% =====================================================================
% Paper 38 (standalone): SU(2)/S^3 state-space Gromov-Hausdorff
% convergence  (the file name retains the superseded 'propinquity'
% wording; Latremoliere propinquity is NOT the distance established)
% on Camporesi-Higuchi truncated metric spectral triples
%
% J. Geom. Phys. style; companion to Leimbach--van Suijlekom 2024
% (Adv. Math. 439, 109496) for flat tori.
%
% Status: first draft, May 2026, post-WH1-PROVEN closure (2026-05-06).
% Built from internal proof memos: r25_l2_proof_memo.md,
% r25_l2_quantitative_rate_memo.md, r25_l3_proof_memo.md,
% r25_l4_proof_memo.md, r25_l5_proof_memo.md, wh1_r35_full_dirac_memo.md.
% =====================================================================
\documentclass[11pt,a4paper]{article}

\usepackage[T1]{fontenc}
\usepackage[utf8]{inputenc}
\usepackage{amsmath, amssymb, amsthm, mathrsfs, mathtools}
\usepackage{geometry}
\geometry{margin=1in}
\usepackage{hyperref}
\hypersetup{colorlinks=true, linkcolor=blue!50!black, citecolor=blue!50!black, urlcolor=blue!50!black}
% \usepackage{microtype}  % disabled: MiKTeX font-expansion environmental issue (same fix as Paper 45)
\usepackage{graphicx}
\usepackage{booktabs}
\usepackage[numbers,sort&compress]{natbib}
\usepackage{xcolor}

\theoremstyle{plain}
\newtheorem{theorem}{Theorem}[section]
\newtheorem{lemma}[theorem]{Lemma}
\newtheorem{proposition}[theorem]{Proposition}
\newtheorem{corollary}[theorem]{Corollary}
\theoremstyle{definition}
\newtheorem{definition}[theorem]{Definition}
\newtheorem{remark}[theorem]{Remark}
\newtheorem{example}[theorem]{Example}

\newcommand{\nmax}{n_{\max}}
\newcommand{\sthree}{S^{3}}
\newcommand{\Hilb}{\mathcal{H}}
\newcommand{\Op}{\mathcal{O}}
\newcommand{\Tcal}{\mathcal{T}}
\newcommand{\Acal}{\mathcal{A}}
\newcommand{\opnorm}[1]{\lVert #1 \rVert_{\mathrm{op}}}
\newcommand{\Lipnorm}[1]{\lVert #1 \rVert_{\mathrm{Lip}}}
\newcommand{\norm}[1]{\lVert #1 \rVert}
\newcommand{\R}{\mathbb{R}}
\newcommand{\C}{\mathbb{C}}
\newcommand{\T}{\mathbb{T}}
\newcommand{\N}{\mathbb{N}}
\newcommand{\Z}{\mathbb{Z}}
\newcommand{\Q}{\mathbb{Q}}
\newcommand{\SU}{\mathrm{SU}}
\newcommand{\SO}{\mathrm{SO}}
\newcommand{\Vol}{\mathrm{Vol}}
\newcommand{\Tr}{\mathrm{Tr}}
\newcommand{\Mat}{M}
\newcommand{\DCH}{D_{\mathrm{CH}}}
\newcommand{\KS}{K}
\newcommand{\Khat}{\widehat{K}}

\title{State-space Gromov--Hausdorff convergence of \\
truncated Camporesi--Higuchi spectral triples on $\sthree$}
\author{J.~Loutey\thanks{Independent researcher. \texttt{jloutey@gmail.com}.
This work was developed in an AI-augmented agentic workflow with
Anthropic's Claude. Implementation, internal memos, and bibliographic
collation were performed collaboratively; all theorems, design choices,
and editorial decisions are the author's responsibility.}}
\date{June 2026 (preprint)}

\begin{document}
\maketitle

\begin{abstract}
\textbf{[INTERNAL THEOREM]} We prove, unconditionally, that the
Connes--van~Suijlekom
truncations of the round $\sthree = \SU(2)$ Camporesi--Higuchi
spectral triple converge to the continuum spectral triple in
van~Suijlekom's state-space Gromov--Hausdorff
distance~\cite{vs2021_jgp}, with explicit rate
\(
d_{\mathrm{GH}} \le \gamma_{\nmax}
= (4/\pi + o(1)) \log \nmax / \nmax
\)
in the dual-Coxeter normalisation of the $\SU(2)$ metric (equivalently $2/\pi$ on the unit round $\sthree$; the constant scales with the metric, 2026-09-03 convention note).
\textbf{[INTERNAL THEOREM]} The truncated state spaces are metrized
by the left-translation
Lipschitz seminorm, which equals the Dirac/metric Lipschitz constant
on the continuum and whose kernel at every finite cutoff is exactly
the scalars (the truthful Dirac-commutator seminorm itself
degenerates at finite cutoff;\ \textbf{[PANEL-VERIFIED]} a remark quantifies
this on a finite panel of cutoffs). The proof
is a compression/lifted-state argument in the pattern of the
compact-metric-group theorem of
Gaudillot-Estrada--van~Suijlekom~\cite{gaudillot_vs2025}, with
self-contained proofs:\ the dual map evaluates conjugated operators
against the central spectral Fej\'er vector state, which embeds
\emph{exactly} into the spinor window via the
$\mathrm{SU}(2)_L \times \mathrm{SU}(2)_R$ block decomposition
$V_j \otimes (V_{j \pm 1/2})$ of the Camporesi--Higuchi eigenspaces,
and both almost-inverse defects collapse to Fej\'er smoothing — one
on functions, one as operator conjugation-averaging — at the common
moment $\gamma_{\nmax}$. The contribution relative to the scalar
theory:\ the spinor-window transport, per-band multiplicity-one
injectivity on the chirality-doubled multiplier system, the
closed-form rate moment, and the $4/\pi$ asymptotic constant.
Supporting structure is organized in five lemmas:\ (L1$'$) a truncated operator
system on the full chirality-doubled spinor bundle;\ (L2)
\textbf{[INTERNAL THEOREM + MEASURED]} a positive
central spectral Fej\'er kernel on $\SU(2)$ with closed-form
mass-concentration moment and Stein--Weiss-type quantitative rate
$\gamma_n \sim (4/\pi)\log n / n$ (dual-Coxeter normalisation; $2/\pi$ on the unit $\sthree$);\ (L3) \textbf{[PANEL-VERIFIED]} a
Lipschitz comparison bound
$\opnorm{[\DCH, \Mat_f]} \leq C_3\,\norm{\nabla f}_{L^\infty}$ with
$C_3 = 1$ (verified on a finite harmonic panel rather than proved in
general, and \emph{not} used by the unconditional theorem);\ (L4) a
Berezin map (convolution form) that is positive,
contractive, an approximate identity, and L3-compatible;\ (L5)
assembly of the distance bound via an approximation pair.
\textbf{[OBSERVATION]} The
asymptotic constant identifies as the Hopf-base measure factor ---
$4/\pi$ in the dual-Coxeter normalisation of \S\ref{sec:ch_triple}
($2/\pi$ on the unit $\sthree$) --- exhibiting the expected one-log slowdown of the
non-abelian SU(2) rate relative to the
Leimbach--van~Suijlekom~\cite{leimbach_vs2024} torus rate $O(1/\nmax)$.
State-space Gromov--Hausdorff convergence of scalar truncations for
\emph{all} compact metric groups — including $\SU(2)$ — is
established by Gaudillot-Estrada--van~Suijlekom~\cite{gaudillot_vs2025};\
the residual contribution here is the chirality-doubled
Camporesi--Higuchi spinor substrate, the closed-form rate moment, and
the explicit $4/\pi$ asymptotic constant. The construction supplies
quantitative metric data for almost-commutative spectral-triple
constructions hosted on $\sthree$~\cite{marcolli_vs2014}.
\end{abstract}

\bigskip\noindent
\textbf{MSC2020 classifications:}\ 58B34 (noncommutative geometry),
46L87 (noncommutative Riemannian and metric geometry),
46L52 (noncommutative function spaces),
22E45 (representations of compact Lie groups).

\medskip\noindent
\textbf{Keywords:}\ spectral triple, state-space Gromov--Hausdorff
convergence, Camporesi--Higuchi Dirac, operator system, Berezin
reconstruction, Peter--Weyl, $\SU(2)$.

% =====================================================================
\section{Introduction}\label{sec:intro}
% =====================================================================

A central problem in noncommutative metric geometry is whether
finite-dimensional truncations of a Riemannian spectral triple
converge, in a suitable quantum Gromov--Hausdorff sense, to the
continuum geometry. Connes and van~Suijlekom~\cite{connes_vs2021}
introduced the framework of \emph{truncated operator systems}
$\Op_\Lambda = P_\Lambda \mathcal{A} P_\Lambda$ obtained by spectral
projection onto the modes of $D$ with $|D| \le \Lambda$, and showed
that the natural propagation invariant $\mathrm{prop}(\Op_\Lambda) = 2$
is realised on the Toeplitz operator system on $S^1$. They left
explicit convergence to the continuum in a quantitative quantum
Gromov--Hausdorff distance as an open question, deferring it to
``elsewhere.'' The convergence itself, at the state-space level, is
established in van~Suijlekom~\cite{vs2021_jgp} for the circle and
sphere, and for \emph{all compact metric groups} by
Gaudillot-Estrada--van~Suijlekom~\cite{gaudillot_vs2025};\ the
residual contribution of the present paper is therefore the spinor
substrate, the closed-form rate moment, and the explicit $4/\pi$
constant — not the existence of convergence for $\SU(2)$ scalars.

The flat-structure case has been worked out in several subsequent
papers. Hekkelman~\cite{hekkelman2022} treated the circle $S^1$
with the Toeplitz/Fej\'er kernel reconstruction. Leimbach and
van~Suijlekom~\cite{leimbach_vs2024}
proved state-space GH convergence of truncated tori with explicit rate
$O(1/\Lambda)$, using the standard
Schur-multiplier characterisation of the abelian-compact-group Fej\'er
mechanism. A complementary line of work treats
\emph{Berezin--Toeplitz} reconstruction on $S^2$, using the K\"ahler
structure of the coadjoint orbit (see~\cite{ucp_maps_2024} and
references therein). All of these results address flat
abelian or K\"ahler structures.

Adjacent activity in the broader propinquity program addresses
orthogonal corners of the same general framework.
Toyota~\cite{toyota2023} establishes quantum-Gromov--Hausdorff
convergence of spectral truncations for discrete groups of
polynomial growth, a scope which excludes compact Lie groups by
definition. Farsi--Latr\'emoli\`ere~\cite{farsi_latremoliere2024}
prove spectral continuity for collapse of spectral triples along
vertical directions in the spectral propinquity, a
collapsing-family setup distinct from the fixed-manifold truncation
sequence here.
Farsi--Latr\'emoli\`ere~\cite{farsi_latremoliere2025} establish
continuity of the spectral propinquity along analytic paths of
Riemannian metrics on a fixed closed spin manifold, a
continuous-family setup distinct from the discrete-truncation
sequence considered here. Hekkelman--McDonald~\cite{hekkelman_mcdonald2024b}
develop a noncommutative integral and a Szeg\"o limit formula for
spectrally truncated triples in the Connes--van~Suijlekom paradigm,
a companion direction complementary to the metric-side convergence
proved here.

For the most natural physically motivated non-K\"ahler compact case
--- the round $\sthree = \SU(2)$ --- the scalar state-space
convergence is covered by the compact-metric-group theorem of
Gaudillot-Estrada--van~Suijlekom~\cite{gaudillot_vs2025};\ what
remains distinctive is the \emph{Dirac spectral triple} version on
the Camporesi--Higuchi spinor substrate, with explicit closed-form
rate.

\textbf{[INTERNAL THEOREM]} The present paper treats that case;\ the
theorem is
unconditional under the translation-seminorm metrization of
\S~\ref{sec:named_gaps}. Let $\Tcal_{\sthree} =
(C^\infty(\sthree), L^2(\sthree, \Sigma), \DCH)$ denote the
\emph{Camporesi--Higuchi spectral triple}, with $\DCH$ the Dirac
operator on the unit round $\sthree$ in the spinor harmonic basis
of Camporesi--Higuchi~\cite{camporesi_higuchi1996}, with spectrum
$\sigma(\DCH) = \{\pm(n + 1/2) : n \in \N\}$ on each chirality
sector.\footnote{We use the Fock-shell labelling $n_{\mathrm{Fock}} =
n_{\mathrm{CH}} + 1$ throughout, so the lowest non-trivial level is
$n=1$ with $|\lambda_1| = 3/2$. This convention is fixed by
\cite[\S2]{camporesi_higuchi1996} and matches the $\SO(4)$
representation-theoretic indexing used in the supporting computational
modules of the present work.} Let $\Tcal_{\nmax} =
(\Op_{\nmax}, \Hilb_{\nmax}, D_{\nmax})$ denote its truncation at
cutoff $\nmax$, where $\Op_{\nmax}$ is the truncated operator system
in the sense of Connes--van~Suijlekom and $D_{\nmax}$ is the
restriction of $\DCH$ to $\Hilb_{\nmax} = P_{\nmax} L^2(\sthree,
\Sigma)$.

\begin{theorem}[Main theorem;\ cf.~\S\ref{sec:named_gaps} and
\S\ref{sec:main_theorem}]
\label{thm:main_intro}
\textbf{[INTERNAL THEOREM]}
Unconditionally, the truncated triples $\Tcal_{\nmax}$ converge to
$\Tcal_{\sthree}$ in van~Suijlekom's state-space Gromov--Hausdorff
distance~\cite{vs2021_jgp}, under the translation-seminorm
metrization of \S~\ref{sec:named_gaps}, with explicit rate
\begin{equation}\label{eq:main_rate}
   \Lambda(\Tcal_{\nmax}, \Tcal_{\sthree})
   \;\le\; C_3 \cdot \gamma_{\nmax}, \qquad
   \gamma_{\nmax} = \frac{4 \log \nmax}{\pi\,\nmax} + O\!\bigl(1/\nmax\bigr)
   \quad\text{(dual-Coxeter normalisation; $2/\pi$ on the unit $\sthree$)},
\end{equation}
where \textbf{[PANEL-VERIFIED]} $C_3 = 1$ is the Lipschitz comparison
constant of
Lemma~\ref{lem:L3} and $\gamma_{\nmax}$ is the mass-concentration
moment of the central spectral Fej\'er kernel on $\SU(2)$ given by
Lemma~\ref{lem:L2}. The unconditional statement is
Theorem~\ref{thm:main_unconditional}, whose bound is
$\le \gamma_{\nmax}$ outright and does not invoke $C_3$;\ the $C_3$
factor displayed here belongs to the Lemma~\ref{lem:L5} assembly route
and carries that lemma's panel scope.
\textbf{[OBSERVATION]} The asymptotic constant is $4/\pi$ in the
dual-Coxeter normalisation of \S\ref{sec:ch_triple} and $2/\pi$ on the
unit $\sthree$;\ it is the sphere-volume quotient $2\Vol(S^2)/\Vol(\sthree)$ for
unit-radius spheres --- though not the Hopf base-to-total ratio (see the
volume-ratio note in \S\ref{sec:L2}) (\emph{metric convention, 2026-09-03:} $\gamma_{\nmax}$ as computed in Lemma~\ref{lem:L2} weights the kernel by the $\SO(3)$ rotation angle $\chi \in [0, 2\pi]$, which is \emph{twice} the geodesic distance $\theta = \chi/2$ on the unit round $\sthree$; the constant is $4/\pi$ in that dual-Coxeter normalisation --- the round $\sthree$ of radius~2 --- and $2/\pi$ on the unit $\sthree$ of \S\ref{sec:setup}, the two differing by the metric scale, which enters a distance moment linearly; the theorem holds in either normalisation, \texttt{tests/test\_p38\_metric\_convention.py}).
\end{theorem}

The proof structure parallels the abelian Leimbach--van~Suijlekom
torus argument but requires substantively more machinery on three
fronts. First, the central spectral Fej\'er kernel on $\SU(2)$ must be
constructed using Peter--Weyl decomposition rather than Schur--Fourier
multipliers; the natural-coefficient version (carrying the
$\sqrt{2j+1}$ Plancherel weight) gives a Ces\`aro-2-style $\log n / n$
rate rather than the unweighted $1/n$ rate of the torus. Second, the
Lipschitz comparison estimate on the round $\sthree$ requires the
Wen--Avery 3-Y integral~\cite{avery_wen_avery1986, paper7} to
control the multiplier matrices in the truncated basis, and depends on
the $\SO(4)$ selection rule
$|n - n'| + 1 \le N \le n + n' - 1$ for non-vanishing 3-Y matrix
elements. Third, the Berezin reconstruction map on $\SU(2)$ does not
arise from a K\"ahler structure (the round $\sthree$ has no compatible
almost-complex structure), so we use the central-Fej\'er-kernel
spectral form directly, which is the natural non-K\"ahler analog of
the Connes--van~Suijlekom $S^1$ Fej\'er reconstruction.

The asymptotic constant in \eqref{eq:main_rate} carries
geometric content, and its value is normalisation-dependent (2026-09-03):\
in the dual-Coxeter normalisation of \S\ref{sec:ch_triple} the constant
is $4/\pi$, which coincides numerically with $\Vol(S^2)/\pi^2$ for
$\Vol(S^2) = 4\pi$;\ on the unit $\sthree$ it is $2/\pi$, the factor $2$
being the metric scale rather than a geometric signature.  \emph{Which volume ratio (2026-09-03).}  The constant is a genuine quotient of sphere volumes --- but not the one its usual name suggests.  For unit-radius spheres $\Vol(S^2)/\Vol(\sthree) = 4\pi/2\pi^2 = 2/\pi$ exactly, and the dual-Coxeter constant is twice that.  It is \emph{not} the Hopf fibration's base-to-total ratio:\ the Hopf map is a Riemannian submersion $\sthree(r) \to S^2(r/2)$, so that ratio is $1/(2\pi)$ on the unit sphere and $1/(4\pi)$ in the dual-Coxeter metric.  The \emph{volume content} of the constant --- what places it in the $k = 0$ (M1) slot of the master Mellin engine and in $\mathbb{Q}[\pi, \pi^{-1}]$ --- is therefore intact;\ only the ``Hopf base'' label is a misnomer, that base carrying half the radius of the sphere whose volume appears.  Accordingly; \textbf{[OBSERVATION]} the same ratio $2/\pi$ is the
first-moment constant of the classical Fej\'er kernel on the circle
(Remark~\ref{rem:circle_fejer}). This is precisely the
\emph{Hopf base measure} factor that appears in the heat-kernel
expansion of the round $\sthree$ via the Hopf fibration $S^1 \to
\sthree \to S^2$, and which has been independently identified in the
internal supporting analysis as one of the three sub-mechanisms of a
master Mellin-engine $\Tr (D^k e^{-tD^2})$ with operator order
$k \in \{0, 1, 2\}$~\cite{paper18}. \textbf{[OBSERVATION]} The constant in
\eqref{eq:main_rate} ($2/\pi$ on the unit $\sthree$; $4/\pi$ in the
dual-Coxeter normalisation) identifies as the $k=0$ sub-case (M1,
Hopf-base measure) of that classification; the same $1/\pi^2$ factor reappears
in the Pythagorean Hilbert--Schmidt orthogonality corollary on the
Krein wedge documented in the companion Lorentzian-extension paper
(Paper~43~\cite{paper43}). \textbf{[OBSERVATION]} The
appearance of this constant in
the state-space GH rate is therefore not an isolated occurrence:\ the
rate at which
the framework recovers its continuum limit carries the same
transcendental signature that the framework's heat-kernel machinery
imprints on its bulk observables.  No derivation links the two
appearances;\ the shared value is the whole of the evidence.

\paragraph{Outline.}
Section~\ref{sec:setup} sets up the round $\sthree$
Camporesi--Higuchi spectral triple, its truncations, and the Peter--Weyl
basis. Section~\ref{sec:lemmas} proves the five lemmas. Each lemma is
self-contained:\ Lemma~\ref{lem:L1prime} treats the chirality-doubled
operator system substrate; Lemma~\ref{lem:L2} the central spectral
Fej\'er kernel and its quantitative rate; Lemma~\ref{lem:L3} the
Lipschitz comparison bound with $C_3 = 1$; Lemma~\ref{lem:L4} the
Berezin reconstruction; Lemma~\ref{lem:L5} the state-space
Gromov--Hausdorff assembly. Section~\ref{sec:main_theorem} assembles the
main theorem. Section~\ref{sec:discussion} discusses the role of the
$4/\pi$ constant, the Marcolli--van~Suijlekom~\cite{marcolli_vs2014}
gauge-network framework that the result anchors, the obstruction to
extending the proof to the holomorphic-sector $S^5$ Bargmann--Segal
case, and open questions.

\subsection{The unconditional theorem:\ compression and lifted
states under the translation seminorm}\label{sec:named_gaps}

Both gaps that conditioned earlier formulations of the main theorem
are closed by one reframing:\ metrize the truncated state space by
the \emph{left-translation Lipschitz seminorm} rather than by the
truthful-Dirac commutator seminorm. On the continuum the two agree
exactly (Lemma~\ref{lem:continuum_lip});\ at finite cutoff the
translation seminorm is non-degenerate while the Dirac-commutator
seminorm is not (Remark~\ref{rem:dirac_degeneracy}). All proofs in
this subsection are self-contained;\ the kernel quantities
$\KS_{\nmax}$, $Z_{\nmax}$, $\gamma_{\nmax}$ are those of
\S~\ref{sec:L2} below (forward references).

\begin{definition}[translation seminorm]
\label{def:translation_seminorm}
Let $\lambda_g$ denote left translation on $\sthree = \SU(2)$, acting
on the spinor space $L^2(\sthree, \Sigma) \cong L^2(\SU(2)) \otimes
\C^2$ (left trivialization) by $U_g = \lambda_g \otimes 1$, and
$\rho_g(T) = U_g T U_g^*$. Since $\DCH$ is built from the
left-invariant frame, $[U_g, \DCH] = 0$;\ hence $U_g$ preserves every
truncation $\Hilb_{\nmax}$ and conjugation preserves the operator
system:\ $\rho_g(P_{\nmax} M_f P_{\nmax}) = P_{\nmax}
M_{\lambda_g f} P_{\nmax}$. For $T \in \Op_{\nmax}$ and
$f \in C(\sthree)$ define
\begin{equation}
L_{\nmax}(T) := \sup_{g \ne e}
\frac{\opnorm{\rho_g(T) - T}}{d(e, g)},
\qquad
L(f) := \sup_{g \ne e}
\frac{\norm{\lambda_g f - f}_{L^\infty}}{d(e, g)},
\end{equation}
with $d$ the dual-Coxeter geodesic distance of
\S\ref{sec:ch_triple} --- equivalently the rotation angle,
$d(e, g) = \chi(g)$ --- which is bi-invariant.
\end{definition}

\begin{lemma}[continuum identification]\label{lem:continuum_lip}
For $f \in C^\infty(\sthree)$,
$L(f) = \mathrm{Lip}_d(f) = \norm{\nabla_d f}_{L^\infty}
= \tfrac{1}{2}\opnorm{[\DCH, M_f]}$, all four quantities taken in the
dual-Coxeter metric $d$ of \S\ref{sec:ch_triple} except the last, in
which $\DCH$ is the unit-metric operator
(Eq.~\eqref{eq:seminorm_normalisation}).
\end{lemma}

\begin{proof}
$\norm{\lambda_g f - f}_\infty = \sup_x |f(g^{-1}x) - f(x)| \le
\mathrm{Lip}(f)\, d(g^{-1}x, x) = \mathrm{Lip}(f)\, d(e, g)$ by
bi-invariance, so $L(f) \le \mathrm{Lip}_d(f)$;\ conversely left
translations act transitively and realize every short geodesic
($d(g^{-1}x, x) = d(e, g)$ along the orbit), so the supremum
recovers $\mathrm{Lip}_d(f)$. The identity $\mathrm{Lip}_d(f) =
\norm{\nabla_d f}_\infty$ on the continuum is standard, as is
$\mathrm{Lip}_{d_1}(f) = \opnorm{[\DCH, M_f]}$ for the unit metric
(Connes' distance formula on a closed spin manifold);\ since
$d = 2 d_1$, Lipschitz seminorms scale inversely and the factor
$\tfrac12$ of \eqref{eq:seminorm_normalisation} follows.
\end{proof}

\begin{lemma}[per-band injectivity]\label{lem:band_injectivity}
For $1 \le N \le 2\nmax - 1$ let $E_N \subset C^\infty(\sthree)$ be
the degree-$(N{-}1)$ harmonic band ($\dim E_N = N^2$), an irreducible
$\mathrm{SO}(4)$-module of type $V_{(N-1)/2} \otimes V_{(N-1)/2}$.
The compression $\iota_N : E_N \to B(\Hilb_{\nmax})$,
$f \mapsto P_{\nmax} M_f P_{\nmax}$, is injective.
\end{lemma}

\begin{proof}
$\iota_N$ is $\mathrm{SO}(4)$-equivariant for the conjugation action
(rotations implemented on $\Hilb_{\nmax}$ commute with $P_{\nmax}$).
Since $E_N$ is irreducible, $\iota_N$ is either zero or injective
(Schur). It is nonzero:\ every multiplier matrix $\Mat_{NLM}$ with
$N \le 2\nmax - 1$ is nonzero, by non-vanishing of the stretched
$3$-$Y$ matrix element within the top kept shell (Wen--Avery
closed form~\cite{avery_wen_avery1986});\ \textbf{[PANEL-VERIFIED]}
verified
symbolically/numerically for every band at $\nmax \le 5$, where the
per-band rank additionally equals the full $N^2$
(\texttt{debug/p38\_g1g2\_band\_diagnostics.py}, frozen in
\texttt{tests/test\_p38\_action\_seminorm.py}).
\end{proof}

\begin{proposition}[kernel condition on the truthful substrate]
\label{prop:kernel_condition}
$L_{\nmax}(T) = 0$ if and only if $T \in \C 1$.
\end{proposition}

\begin{proof}
$L_{\nmax}(T) = 0$ means $\rho_g(T) = T$ for all $g$. Write $T =
P_{\nmax} M_f P_{\nmax}$ with $f$ band-limited to $N \le 2\nmax - 1$
(every element of $\Op_{\nmax}$ has this form). Then $P_{\nmax}
M_{\lambda_g f - f} P_{\nmax} = 0$;\ left translation preserves each
harmonic band, so $\lambda_g f - f$ is band-limited and
Lemma~\ref{lem:band_injectivity} gives $\lambda_g f = f$ for all
$g$;\ transitivity forces $f$ constant. The converse is immediate.
Finite dimensionality plus the kernel condition make
$(S(\Op_{\nmax}), \mathrm{MK}_{L_{\nmax}})$ a compact metric space of
finite diameter, and Lemma~\ref{lem:continuum_lip} provides the same
on the continuum side --- the standard Monge--Kantorovich setting on
both ends.
\end{proof}

\begin{lemma}[spinor window inclusion and the lifted state]
\label{lem:lifted_state}
Let $J = (\nmax - 1)/2$, let
$h = \sum_{j \le J} \sqrt{2j+1}\,\chi_j$ be the central spectral
Fej\'er amplitude of Definition~\ref{def:central_fejer} (the scalar
window then has exactly $\nmax$ levels, and $|h|^2 / Z_{\nmax}$ is
the Fej\'er kernel $\KS_{\nmax}$), fix a unit spinor $\chi \in
\C^2$, and set $\xi := (h \otimes \chi)/\sqrt{Z_{\nmax}}$. Then:
\begin{enumerate}\setlength{\itemsep}{2pt}
\item[(a)] (\emph{Exact window inclusion.}) $\xi \in \Hilb_{\nmax}$.
Under $\mathrm{SU}(2)_L \times \mathrm{SU}(2)_R$,
\begin{equation*}
L^2(\SU(2)) \otimes \C^2 \;=\; \bigoplus_{j} V_j \otimes
\bigl(V_{j+1/2} \oplus V_{j-1/2}\bigr),
\end{equation*}
and these summands are precisely the Camporesi--Higuchi
eigenspaces~\cite{camporesi_higuchi1996}:\ in the level index of
\eqref{eq:CH_spectrum}, $V_j \otimes V_{j+1/2}$ is the positive shell
$n = 2j + 1$ and $V_j \otimes V_{j-1/2}$ the negative shell $n = 2j$,
with dimensions $n(n+1)$ matching shell by shell. Hence $h \otimes
\chi$ lies in the union of the \emph{complete} shells $n \le 2J + 1 =
\nmax$ --- exact containment, no spillover.
\item[(b)] (\emph{Lifted state.}) For every $f \in C(\sthree)$,
$\langle \xi, (P_{\nmax} M_f P_{\nmax})\,\xi\rangle =
\int_{\sthree} f\, \KS_{\nmax}\, d\mu_{\mathrm{Haar}}$:\ the induced
measure is exactly the Fej\'er measure, with first moment
$\gamma_{\nmax}$ (Lemma~\ref{lem:L2}(d)).
\item[(c)] (\emph{Dual map.}) $\upsilon : \Op_{\nmax} \to
C(\sthree)$, $\upsilon(T)(g) := \langle U_g \xi,\, T\, U_g
\xi\rangle$, is unital and positive, satisfies the tautological
contraction $L(\upsilon(T)) \le L_{\nmax}(T)$, and acts on compressed
multipliers as Fej\'er smoothing:\
$\upsilon(P_{\nmax} M_f P_{\nmax}) = \KS_{\nmax} * f$.
\end{enumerate}
\end{lemma}

\begin{proof}
(a) is the displayed decomposition ($\C^2$ in the right tensor slot;\
$V_j \otimes V_{1/2} = V_{j+1/2} \oplus V_{j-1/2}$), with the
eigenvalue and dimension identifications of
Camporesi--Higuchi~\cite{camporesi_higuchi1996}. (b):\ $\xi \in
\mathrm{ran}\, P_{\nmax}$ absorbs the projections, so the pairing is
$\int f\, |h|^2\, |\chi|^2 / Z_{\nmax} = \int f\, \KS_{\nmax}$.
(c):\ positivity and unitality are immediate from the vector-state
form. The contraction:\ for $q = g' g^{-1}$,
$|\upsilon(T)(g) - \upsilon(T)(g')| = |\langle \xi,
(\rho_{g^{-1}}(T) - \rho_{g'^{-1}}(T))\, \xi\rangle| \le
\opnorm{\rho_q(T) - T} \le L_{\nmax}(T)\, d(e, q) =
L_{\nmax}(T)\, d(g, g')$ by bi-invariance. The smoothing identity:\
$\upsilon(P M_f P)(g) = \int f(x)\, \KS_{\nmax}(g^{-1}x)\, dx =
(\KS_{\nmax} * f)(g)$, using centrality and inversion-invariance of
$\KS_{\nmax}$;\ \textbf{[MEASURED]} verified to machine precision
($\le 3 \times
10^{-16}$ across bands and random group elements,
\texttt{debug/p38\_g1g2\_scalar\_prototype.py}).
\end{proof}

\begin{theorem}[Main theorem, unconditional;\ \textbf{[INTERNAL THEOREM]}]
\label{thm:main_unconditional}
For every $\nmax \ge 1$,
\begin{equation}\label{eq:unconditional_bound}
d_{\mathrm{GH}}\Bigl( \bigl(S(\Op_{\nmax}),
\mathrm{MK}_{L_{\nmax}}\bigr),\;
\bigl(S(C(\sthree)), \mathrm{MK}_{L}\bigr) \Bigr)
\;\le\; \gamma_{\nmax}
\;=\; \frac{4}{\pi}\,\frac{\log \nmax}{\nmax} + O(1/\nmax)
\end{equation}
in the dual-Coxeter normalisation of Lemma~\ref{lem:L2}, i.e.\
$\gamma_{\nmax} = 2\int \KS_{\nmax}\, d_{\mathrm{round}}\, dg$ on the
unit $\sthree$, where the same bound reads $(2/\pi)\log\nmax/\nmax$
(2026-09-03 convention note; the theorem holds in either normalisation).
\end{theorem}

\begin{proof}
Let $S : C(\sthree) \to \Op_{\nmax}$, $S(f) = P_{\nmax} M_f
P_{\nmax}$ (unital completely positive) and let $\upsilon$ be the
dual map of Lemma~\ref{lem:lifted_state}. Four facts:
\begin{enumerate}\setlength{\itemsep}{1pt}
\item[(i)] $L_{\nmax}(S(f)) \le L(f)$:\ $\rho_g(S(f)) - S(f) =
S(\lambda_g f - f)$ and compression contracts operator norms.
\item[(ii)] $L(\upsilon(T)) \le L_{\nmax}(T)$
(Lemma~\ref{lem:lifted_state}(c)).
\item[(iii)] $\norm{\upsilon(S(f)) - f}_\infty =
\norm{\KS_{\nmax} * f - f}_\infty \le \gamma_{\nmax}\, L(f)$ (the
good-kernel estimate in the proof of Lemma~\ref{lem:L4}(c), at the
function level).
\item[(iv)] $\opnorm{S(\upsilon(T)) - T} \le \gamma_{\nmax}\,
L_{\nmax}(T)$:\ on $T = P M_f P$,
\begin{equation*}
S(\upsilon(T)) = P\, M_{\KS_{\nmax} * f}\, P
= \int_{\SU(2)} \rho_g(T)\, \KS_{\nmax}(g)\, dg =: \Phi(T)
\end{equation*}
--- the operator-side conjugation average against the same Fej\'er
measure --- and $\opnorm{\Phi(T) - T} \le \int \opnorm{\rho_g(T) -
T}\, \KS_{\nmax}(g)\, dg \le L_{\nmax}(T) \int d(e, g)\,
\KS_{\nmax}(g)\, dg = \gamma_{\nmax}\, L_{\nmax}(T)$.
\end{enumerate}
By (i)--(ii) the state-space pullbacks $S^* : S(\Op_{\nmax}) \to
S(C(\sthree))$ and $\upsilon^* : S(C(\sthree)) \to S(\Op_{\nmax})$
are $1$-Lipschitz for the Monge--Kantorovich metrics;\ by
(iii)--(iv) they are $\gamma_{\nmax}$-almost mutually inverse, e.g.\
$d(S^*\upsilon^*\omega, \omega) = \sup_{L(f) \le 1}
|\omega(\upsilon(S(f))) - \omega(f)| \le \gamma_{\nmax}$. The
correspondence $\{(\omega, \upsilon^*\omega)\} \cup \{(S^*\phi,
\phi)\}$ then has distortion $\le 2\gamma_{\nmax}$ (the standard
four-pair approximate-isometry computation), whence
$d_{\mathrm{GH}} \le \gamma_{\nmax}$. The rate is
Lemma~\ref{lem:L2}(d.ii)--(d.iii). \qedhere
\end{proof}

\begin{remark}[no inverse estimate;\ the Berezin pair as refinement]
\label{rem:no_inverse}
Both almost-inverse defects in the proof are Fej\'er smoothings at
the same moment $\gamma_{\nmax}$ --- one on functions, one as the
conjugation average $\Phi$ --- so no partial-inverse or transference
estimate is required anywhere:\ the dual-reach gap of earlier
formulations is dissolved rather than closed. The Berezin map
$B_{\nmax}$ of Definition~\ref{def:berezin} (Lemmas~L4--L5) remains
as an explicit forward-reconstruction refinement;\ it is no longer
load-bearing for the distance bound. The scalar $C(G)$ case of this
compression/lifted-state pattern, for general compact metric groups
with the same single-moment bound, is due to
Gaudillot-Estrada--van~Suijlekom~\cite{gaudillot_vs2025};\ the
contribution here is the spinor-window transport (the exact-fit
inclusion of Lemma~\ref{lem:lifted_state}(a)), per-band
multiplicity-one injectivity on the chirality-doubled system, the
closed-form rate moment, and the $4/\pi$ asymptotic constant. All
assembly steps above are proved self-contained.
\end{remark}

\begin{remark}[degeneracy of the Dirac-commutator seminorm]
\label{rem:dirac_degeneracy}
\textbf{[MEASURED]} The truthful chirality-diagonal commutator seminorm
$\opnorm{[\DCH, \cdot]}$ on $\Op_{\nmax}$ has kernel strictly larger
than the scalars at the cutoffs measured ($10/14$ multipliers at
$\nmax = 2$, $26/55$ at $\nmax = 3$, bit-exact;\ every top-band
multiplier compresses shell-diagonally), so it does not metrize the
truncated state space.  (The frozen falsifier exercises the
$\nmax = 2$ counts;\ the $\nmax = 3$ counts are measured but not
frozen.) An engineered off-diagonal modification of
$\DCH$ restores its kernel condition ($1/14$, $1/55$) but is no
longer needed for the main theorem. Frozen falsifier for the kernel
counts:\
\texttt{tests/test\_p45\_kplus\_degeneracy.py::test\_spatial\_kernel\_condition\_truthful\_vs\_offdiag}
(\texttt{tests/test\_p38\_action\_seminorm.py} backs the injectivity and
smoothing premises, not these counts; pointer corrected 2026-09-03).
\end{remark}

\begin{remark}[history]\label{rem:history38}
An earlier draft of this paper (May 2026;\ repository history and the
corresponding Zenodo deposit) stated the main theorem in the
Latr\'emoli\`ere propinquity (a formulation since withdrawn), with a Lemma~L2 that conflated the
kernel's Plancherel mass distribution with its multiplier symbol and
a dual-reach argument resting on that conflation. A verification
pass (2026-06-09;\ repository \texttt{CHANGELOG}, v3.106.0) produced
an intermediate conditional form;\ the present unconditional form
(2026-06-10) follows from the translation-seminorm reframing of this
subsection. The companion Lorentzian paper proves the related
$K^{+}$-compression degeneracy theorem.
\end{remark}

% =====================================================================
\section{Setup}\label{sec:setup}
% =====================================================================

\paragraph{Metric convention.}
The distance in which convergence is proved in this paper is
van~Suijlekom's \emph{state-space Gromov--Hausdorff
distance}~\cite{vs2021_jgp}:\ whenever we write
$\Lambda(\Tcal_1, \Tcal_2)$ we mean
$\mathrm{d}_{\mathrm{GH}}^{\,\mathrm{vS}}(S(\Op_1), S(\Op_2))$, the
Gromov--Hausdorff distance between the state spaces of the associated
operator systems in their Monge--Kantorovich metrics;\ the definition
is restated as \eqref{eq:propinquity_def} in \S\ref{sec:L5}. The
strictly stronger \emph{Gromov--Hausdorff propinquity for metric
spectral triples} of
Latr\'emoli\`ere~\cite{latremoliere_metric_st_2017}
(Adv.~Math. \textbf{404} (2022), Paper No.~108393, preprint
\texttt{arXiv:1811.10843}, 2017) --- the spectral-triple-level
extension of the dual Gromov--Hausdorff propinquity of
\cite{latremoliere2015} (J.~Math.~Pures Appl. \textbf{103} (2015))
and of the original Gromov--Hausdorff propinquity for quantum compact
metric spaces of \cite{latremoliere2016} (Trans.~AMS \textbf{368}
(2016)) --- is \emph{not} the distance established here:\
\textbf{[OPEN]} its tunnel
hypotheses are not verified for the truncated operator system
(\S\ref{sec:L5}), and convergence in that propinquity remains the
strictly stronger open target.

\subsection{The round $\sthree$ Camporesi--Higuchi spectral triple}
\label{sec:ch_triple}

Take $G = \SU(2)$ with the bi-invariant metric in the
\emph{dual-Coxeter normalisation}:\ the $\mathrm{Ad}$-invariant inner
product on $\mathfrak{su}(2)$ fixed by $\mathrm{Cas}(\mathrm{ad}) =
h^\vee = 2$, namely $\langle X, Y\rangle = -2\,\mathrm{tr}(XY)$ in the
defining representation.  On $\SU(2)$ this normalisation makes the
geodesic distance from the identity equal to the \emph{rotation angle}:
for $g = \exp(i\theta\,\hat n \cdot \sigma)$, which has eigenphases
$\pm\theta$, one computes $\lVert i\theta\,\hat n \cdot \sigma\rVert =
2\theta = \chi(g)$.  The underlying manifold is thus the round 3-sphere
of radius~$2$, with diameter $2\pi$ and total volume
$\Vol(\sthree) = 16\pi^2$;\ the \emph{unit} round $\sthree$ is the same
manifold with every distance halved, and we write $d_1 = d/2$ and
$\Vol_1(\sthree) = 2\pi^2$ when that comparison is needed.

We fix the metric by this representation-theoretic rule, rather than by
a choice of radius, because it is the rule under which the rate
constant of Theorem~\ref{thm:main} is comparable across groups:\
Paper~40~\cite{paper40_unified} uses the same normalisation for the
whole compact connected semisimple class, so the present paper is
exactly its rank-1 case.  (Until 2026-09-03 this subsection declared
the unit sphere while the moment $\gamma_{\nmax}$ of
Lemma~\ref{lem:L2} was computed against $\chi$;\ the two differ by the
factor $2$ above, and the bound was correspondingly stated with a
unit-metric seminorm against a dual-Coxeter moment.  Declaring the
metric uniformly removes that seam;\ since every quantity in the bound
is homogeneous of degree one in the metric, the inequality itself is
unchanged.)  Identify $\sthree$
with the Lie group $\SU(2)$ via the standard double cover\footnote{We
use this identification implicitly throughout. It is convenient
because (i)~the Peter--Weyl decomposition of $L^2(\SU(2))$ provides a
natural orthonormal basis indexed by $j \in \tfrac{1}{2} \N_{\ge 0}$,
(ii)~the central conjugacy-class measure on $\SU(2)$ underlies the
mass-concentration analysis of Lemma~\ref{lem:L2}, and (iii)~the
spectrum of $\DCH$ has the simplest closed form in this convention.
We do not draw a structural distinction between $\sthree$ and
$\SU(2)$ in this paper.}. Let $\Sigma$ denote the spinor bundle and
$L^2(\sthree, \Sigma)$ the spinor Hilbert space.

The Camporesi--Higuchi Dirac operator
$\DCH: \mathrm{Dom}(\DCH) \subset L^2(\sthree, \Sigma) \to L^2(\sthree,
\Sigma)$ is the Dirac operator of the \emph{unit} round metric --- we
keep that normalisation of the operator throughout, so that
\eqref{eq:CH_spectrum} below is the familiar Camporesi--Higuchi
spectrum;\ the Dirac operator of the dual-Coxeter metric declared above
is $\tfrac{1}{2}\DCH$.  It acts on the spinor harmonics
\begin{equation}\label{eq:CH_spectrum}
   \DCH \, \psi_{n, l, m_j, \chi}
   \;=\; \chi\,(n + \tfrac{1}{2})\, \psi_{n, l, m_j, \chi},
\end{equation}
with $n \in \N_{\ge 1}$, $l \in \{0, 1, \ldots, n-1\}$, $j = l + 1/2$,
$m_j \in \{-j, -j+1, \ldots, j\}$, and chirality $\chi \in \{+1, -1\}$
\cite[\S 2]{camporesi_higuchi1996}. The degeneracy of the level
$n$ on each chirality sector is $n(n+1)$; the total degeneracy
including both chiralities is $2n(n+1)$.

We then form the Camporesi--Higuchi spectral triple
\begin{equation}
\Tcal_{\sthree}
\;=\; \bigl(\, C^\infty(\sthree),\; L^2(\sthree, \Sigma),\; \DCH \,\bigr).
\end{equation}
This is a unital, even, real spectral triple of metric dimension
$3$ in the sense of Connes~\cite{connes1995}, with KO-dimension $3 \pmod
8$. The chirality grading $\gamma_5 = \mathrm{diag}(+1, -1)$ on the
spinor bundle, the unitary real structure $J$ given by charge
conjugation, and the Lipschitz seminorm all enter as expected.  Because
$\DCH$ is the unit-metric Dirac operator while $d$ is the dual-Coxeter
distance, the two differ by the same factor $2$:
\begin{equation}\label{eq:seminorm_normalisation}
   L(f) \;=\; \mathrm{Lip}_d(f) \;=\; \tfrac{1}{2}\,\opnorm{[\DCH, \Mat_f]},
\end{equation}
and it is $L$ --- the dual-Coxeter Lipschitz seminorm, equivalently the
translation seminorm of Definition~\ref{def:translation_seminorm} ---
that metrizes every state space in this paper.

\subsection{Truncations}\label{sec:truncations}

Let $P_{\nmax}: L^2(\sthree, \Sigma) \to \Hilb_{\nmax}$ denote the
orthogonal projection onto the finite-dimensional subspace spanned by
$\{\psi_{n, l, m_j, \chi} : 1 \le n \le \nmax\}$. The truncated
Hilbert space has dimension
\begin{equation}
   \dim \Hilb_{\nmax}
   \;=\; 2 \sum_{n=1}^{\nmax} n(n+1)
   \;=\; \tfrac{2}{3}\, \nmax(\nmax+1)(\nmax+2).
\end{equation}
The truncated Dirac operator is $D_{\nmax} = P_{\nmax} \DCH P_{\nmax}$,
which (since $\DCH$ is diagonal in the basis above) acts as
$\chi (n+1/2)$ on each kept basis vector.

Following Connes--van~Suijlekom~\cite[\S 3]{connes_vs2021}, the
\emph{truncated operator system} is
\begin{equation}\label{eq:O_def}
   \Op_{\nmax}
   \;=\; P_{\nmax}\, C^\infty(\sthree)\, P_{\nmax}
   \;\subset\; B(\Hilb_{\nmax}),
\end{equation}
viewed as a self-adjoint subspace of the matrix algebra
$\Mat_{\dim \Hilb_{\nmax}}(\C)$. As emphasised in
\cite[\S 1]{connes_vs2021}, this is \emph{not} a $C^*$-subalgebra of
$\Mat_{\dim \Hilb_{\nmax}}(\C)$:\ the full $C^*$-algebra it generates
is the full matrix algebra (the ``$C^*$-envelope''), but the operator
system $\Op_{\nmax}$ itself encodes finer structure, including the
algebra-side analog of the Lipschitz seminorm. The truncated metric
spectral triple is then
\begin{equation}
   \Tcal_{\nmax} \;=\; (\Op_{\nmax}, \Hilb_{\nmax}, D_{\nmax}).
\end{equation}

\subsection{Peter--Weyl basis and the Avery 3-Y integral}
\label{sec:peter_weyl}

The Peter--Weyl decomposition of $L^2(\SU(2))$ gives an orthonormal
basis $\{D^j_{m, m'} : j \in \tfrac{1}{2}\N_{\ge 0}, |m|, |m'| \le j\}$
of Wigner $D$-matrix elements. Setting $n = 2j + 1$ identifies the
Peter--Weyl level with the Fock-shell index used above. We work in the
\emph{Wen--Avery hyperspherical harmonic} basis
\begin{equation}\label{eq:avery_y3}
   Y^{(3)}_{NLM}(\chi, \theta, \phi)
   \;=\; \mathcal{N}_{NL}\,\sin^L \chi \cdot
   C^{L+1}_{N - L - 1}(\cos \chi)\,Y_{LM}(\theta, \phi),
\end{equation}
which is the standard $\SO(4)$-adapted basis of $L^2(\sthree)$
in the form given by Wen and Avery~\cite{avery_wen_avery1986} and used systematically
in the molecular orbital literature; here
$\mathcal{N}_{NL} = \sqrt{2^{2L+1} N (N-L-1)! (L!)^2 / (N+L)!}/\sqrt{\pi}$
is the standard normalisation, $C^{L+1}_{N-L-1}$ is the Gegenbauer
polynomial, and $Y_{LM}$ are the standard spherical harmonics on
$S^2$. The Laplace--Beltrami spectrum on $\sthree$ is then
$-(N^2 - 1)$ on the isotypic component spanned by
$\{Y^{(3)}_{NLM}\}_{L,M}$.

For $f = \sum_{N \ge 1} \sum_{L=0}^{N-1} \sum_{|M| \le L}
c_{NLM}\, Y^{(3)}_{NLM}$, the multiplier matrix
$\Mat_f := P_{\nmax} f P_{\nmax}$ acts on $\Hilb_{\nmax}$ with entries
\begin{equation}\label{eq:multiplier_3y}
   (\Mat_f)_{(n', l', m'_j),\,(n, l, m_j)}
   \;=\; \sum_{N, L, M} c_{NLM}\,
   \mathcal{G}^{(3)}_{NLM; n l m_j; n' l' m'_j},
\end{equation}
where the Avery 3-Y integral $\mathcal{G}^{(3)}$ is the Clebsch--Gordan-
weighted sum of Gegenbauer triple integral and $S^2$ Gaunt integral
factors\footnote{In SO(4) the relevant 3-Y integral on $\sthree$
factors as a Gegenbauer triple integral times a standard 3-Y
integral on $S^2$, with the $\SO(4)$ selection rule
$|n-n'|+1 \le N \le n+n'-1$ enforcing the truncation
$|n-n'| < N$ on the radial channel and the
standard $S^2$ Gaunt selection rule on the angular channel. This
double-selection structure is essential to Lemma~\ref{lem:L3} and is
treated as canonical infrastructure here. Closed-form values used in
internal numerical verifications appear in the supporting code module
\texttt{geovac/so4\_three\_y\_integral.py} (R3.1, May 2026).}.

The selection rule $|n - n'| + 1 \le N \le n + n' - 1$ enforces a
sparse structure on $\Mat_f$ that we use crucially in
Lemma~\ref{lem:L3}.

% =====================================================================
\section{The five lemmas}\label{sec:lemmas}
% =====================================================================

% ---------------------------------------------------------------------
\subsection{Lemma L1$'$: chirality-doubled operator system}
\label{sec:L1prime}
% ---------------------------------------------------------------------

The chirality-doubled spinor bundle $\Sigma = \Sigma_+ \oplus \Sigma_-$
gives the truncated Hilbert space the structure
$\Hilb_{\nmax} = \Hilb^+_{\nmax} \oplus \Hilb^-_{\nmax}$, with each
chirality block carrying a Weyl-sector multiplier matrix
$\Mat_f^{\mathrm{Weyl}}$ acting identically on the two blocks. The
multiplier $\Mat_f$ is therefore block-diagonal in chirality:
\begin{equation}\label{eq:multiplier_block_diagonal}
   \Mat_f \;=\; \Mat_f^{\mathrm{Weyl}} \oplus \Mat_f^{\mathrm{Weyl}}
   \quad\text{on}\quad \Hilb^+_{\nmax} \oplus \Hilb^-_{\nmax}.
\end{equation}

\begin{lemma}[L1$'$:\ chirality-doubled operator system substrate;\
\textbf{[SYMBOLIC + MEASURED]}]
\label{lem:L1prime}
The operator system $\Op_{\nmax}$ inherits from
\eqref{eq:multiplier_block_diagonal} a chirality block-diagonal
structure. The truncated Connes distance
\begin{equation}\label{eq:connes_distance_def}
   d_{\Tcal_{\nmax}}(\varphi, \psi)
   \;=\; \sup\bigl\{\,|\varphi(M) - \psi(M)| \;:\;
   M \in \Op_{\nmax},\, \opnorm{[D_{\nmax}, M]} \le 1 \,\bigr\}
\end{equation}
is finite on every cross-shell pair of pure node-evaluation states
when \eqref{eq:multiplier_block_diagonal} is replaced by the
\emph{chirality-mixing offdiagonal Camporesi--Higuchi extension}
$D_{\nmax}^{\mathrm{offdiag}}$ obtained by adding the standard
$E1$ ladder couplings between $(n, l, m_j, +)$ and $(n+1, l \pm 1, m_j,
-)$.
\end{lemma}

\begin{proof}
The block-diagonal structure of $\Mat_f$ in \eqref{eq:multiplier_block_diagonal}
follows directly from the chirality-blindness of scalar multipliers. To
verify finiteness of the Connes distance on cross-shell pairs, we
observe that a state pair $(\varphi_{n, l, m_j}, \psi_{n', l', m'_j})$
with $n \ne n'$ has finite truncated Connes distance under the offdiag
extension because the $E1$ coupling supplies a non-trivial structural
multiplier that bounds the supremum in
\eqref{eq:connes_distance_def}. \textbf{[PANEL-VERIFIED]} Numerical
verification at
$\nmax \in \{2, 3\}$ shows every cross-shell pair has finite distance,
with the four-point Pearson correlation between distance and Fock-graph
hop distance taking the value $-0.2501$ at $\nmax = 3$;\ no frozen
test in the repository backs this correlation value.\footnote{The
numerical correlation value is recorded in the internal computation
log; the structural finiteness is the load-bearing input for
Lemma~\ref{lem:L5}, and is robust under operator-norm perturbations of
$D$ that preserve the SU(2) selection rules.}

The role of the offdiag extension is purely
SDP-bounding:\ it ensures that the operator-norm constraint
$\opnorm{[D_{\nmax}, M]} \le 1$ is non-degenerate on the full operator
system. The truthful $\DCH$ (without offdiag couplings) is what
underlies the spectral-action structure of
$\Tcal_{\sthree}$ and is what we use in Lemmas~\ref{lem:L3}--\ref{lem:L4}.
The two operators are related by a bounded compact perturbation, and
the state-space GH bound of Theorem~\ref{thm:main_intro} is robust under
this choice.
\end{proof}

\begin{remark}\label{rem:offdiag}
The offdiag extension does \emph{not} satisfy
$\{\gamma_5, D^{\mathrm{offdiag}}_{\nmax}\} = 0$, whereas the truthful
$\DCH$ does. The offdiag is therefore a metric-only device for
finite-distance verification, not a candidate for spectral-action
analysis. The truthful $\DCH$ is what enters
\eqref{eq:CH_spectrum} and what we use throughout the state-space GH
convergence proof.
\end{remark}

% ---------------------------------------------------------------------
\subsection{Lemma L2:\ central spectral Fej\'er kernel on SU(2)}
\label{sec:L2}
% ---------------------------------------------------------------------

We construct the central spectral Fej\'er kernel
$\KS_{\nmax} \in C(\SU(2))$ that plays the role of the
abelian-compact-group Fej\'er kernel in Connes--van~Suijlekom
\cite[\S 3]{connes_vs2021} but with the natural-coefficient Plancherel
weights forced by the Peter--Weyl bijection $n = 2j + 1$.

\begin{definition}[Central spectral Fej\'er kernel]\label{def:central_fejer}
For $\nmax \ge 1$, define
\begin{equation}\label{eq:K_def}
   \KS_{\nmax}(g)
   \;:=\; \frac{1}{Z_{\nmax}}\,
   \biggl|\,\sum_{j \le j_{\max}} \sqrt{2j+1}\, \chi_j(g)\,\biggr|^{2},
   \qquad Z_{\nmax} = \frac{\nmax(\nmax+1)}{2},
\end{equation}
where $\chi_j$ is the trace character of the $(2j+1)$-dimensional
irreducible representation of $\SU(2)$ and $j_{\max}$ is the largest
half-integer with $2j_{\max}+1 \le \nmax$.
\end{definition}

The kernel is positive on $\SU(2)$ (manifestly), central (depends only
on the conjugacy class of $g$), and unit-mass under the standard
$\SU(2)$ Haar measure $dg$. Identifying conjugacy classes with the
rotation angle $\chi \in [0, 2\pi]$ via $g = \exp(i (\chi/2)\, \hat{n}
\cdot \sigma)$ on the maximal torus (eigenphases $\pm\chi/2$;
unit-$\sthree$ geodesic distance $\chi/2$; corrected 2026-09-03 from
$\chi \in [0, \pi]$ with $g = \exp(i\chi \hat n \cdot \sigma)$, which
is the geodesic angle and does not match the characters below), the
conjugacy-class measure is $d\mu(\chi) = (1/\pi) \sin^2(\chi/2)\, d\chi$.

\begin{lemma}[L2: kernel construction and quantitative rate]\label{lem:L2}
The central spectral Fej\'er kernel $\KS_{\nmax}$ has the following
properties (the closed-form $\gamma_{\nmax}$ rate moment is verified by
\texttt{tests/test\_central\_fejer\_su2.py}, and the $4/\pi$ asymptotic
constant by \texttt{tests/test\_trunk\_qa\_fejer\_4\_over\_pi.py};\ $2/\pi$
is the same constant on the unit $\sthree$ and the circle's constant,
not a decoy):
\begin{enumerate}\setlength{\itemsep}{2pt}
\item[(a)] $\KS_{\nmax}(g) \ge 0$ for all $g$, and $\int_{\SU(2)}
\KS_{\nmax}(g) \, dg = 1$.
\item[(b)] (Mass distribution and multiplier symbol.) The kernel's Plancherel \emph{mass distribution} is
$m_{\nmax}(j) = (2j+1)/Z_{\nmax}$ on $j \le j_{\max}$ (the squared
Plancherel weights of the amplitude
$\sum_{j}\sqrt{2j+1}\,\chi_{j}$), a probability distribution with
maximum $2/(\nmax+1)$ at $j_{\max}$. The convolution operator
$T_{\KS_{\nmax}} : f \mapsto \KS_{\nmax} \ast f$ acts on the
$j'$-isotypic component by the \emph{multiplier symbol}
$\sigma_{\nmax}(j') = c_{j'} / ((2j'+1) Z_{\nmax})$ with
$c_{j'} = \sum_{j_1, j_2 \le j_{\max},\, \triangle(j_1, j_2, j')}
\sqrt{(2j_1+1)(2j_2+1)}$, where $\triangle$ is the Clebsch--Gordan
condition \emph{with unit steps} ($|j_1 - j_2| \le j' \le j_1 + j_2$
and $j_1 + j_2 + j' \in \mathbb{Z}$):\
$\sigma_{\nmax}(0) = 1$, $0 \le \sigma_{\nmax} \le 1$, support
$j' \le 2j_{\max}$. \textbf{[MEASURED]} At $\nmax = 2$:\
$\sigma = (1, \sqrt2/3, 2/9)$ at $j' = (0, \tfrac12, 1)$
(bit-exact against independent Weyl-integration quadrature).
\item[(c)] (Cb-norm.) $T_{\KS_{\nmax}}$
is unital completely positive (convolution against a probability
density), hence $\Vert T_{\KS_{\nmax}}\Vert_{\mathrm{cb}} = 1$. (The
mass-distribution maximum $2/(\nmax+1)$ of (b) is \emph{not} the
cb-norm of any map used in this paper.)
\item[(d)] (Mass-concentration moment.) Define
\begin{equation}\label{eq:gamma_def}
   \gamma_{\nmax}
   \;:=\; \int_{\SU(2)} \KS_{\nmax}(g)\, \chi(g)\, dg,
\end{equation}
where $\chi(g) \in [0, 2\pi]$ is the class (rotation) angle of $g$,
which by \S\ref{sec:ch_triple} \emph{is} the dual-Coxeter geodesic
distance $d(e, g)$ ($g = \cos\theta + i \sin\theta\, \hat n \cdot
\sigma$ has eigenphases $\pm\theta$ and $d(e,g) = 2\theta = \chi$;\ the
characters below are $\sin((2j{+}1)\chi/2)/\sin(\chi/2)$).  So
$\gamma_{\nmax}$ is the mass-concentration moment against the paper's
own metric, and every constant below is a dual-Coxeter constant;\ the
corresponding unit-$\sthree$ moment is exactly $\gamma_{\nmax}/2$ at
every $\nmax$ (\texttt{tests/test\_p38\_metric\_convention.py};\
$\KS_1 \equiv 1$ makes $\gamma_1 = \pi$, exactly twice the unit-$\sthree$
Haar-mean distance $\pi/2$).  \emph{Correction (2026-09-03):} this
definition previously wrote ``$d_{\mathrm{round}}$'' while
\S\ref{sec:ch_triple} declared the unit sphere, so the moment and the
seminorm were stated in different metrics;\ both are now dual-Coxeter.
Then:
\begin{enumerate}\setlength{\itemsep}{1pt}
\item[(d.i)] \textbf{[INTERNAL THEOREM]} $\gamma_{\nmax}$ admits the
closed-form sum-rule
\begin{equation}\label{eq:gamma_sum_rule}
   \gamma_{\nmax}
   \;=\; \pi \;-\; \frac{4 T_{\nmax}}{\pi\,Z_{\nmax}}, \qquad
   T_{\nmax}
   \;:=\!\!\!\sum_{\substack{1\le k_1, k_2 \le \nmax \\ k_1 + k_2\, \mathrm{odd}}}
   \sqrt{k_1 k_2}\,\Big[\frac{1}{(k_1 - k_2)^2}
                       - \frac{1}{(k_1 + k_2)^2}\Big].
\end{equation}
\item[(d.ii)] (Asymptotic rate.) \textbf{[INTERNAL THEOREM +
MEASURED]} $\displaystyle
\lim_{\nmax \to \infty} \frac{\nmax\,\gamma_{\nmax}}{\log \nmax}
\;=\; \frac{4}{\pi}.
$
\item[(d.iii)] (Uniform bound.) \textbf{[PANEL-VERIFIED]}
$\gamma_{\nmax} \le 6\,\log \nmax / \nmax$
for all $\nmax \ge 2$ (established through a monotonicity statement
checked on a finite range of $\nmax$, not proved for all $\nmax$).
\end{enumerate}
\end{enumerate}
\end{lemma}

\begin{proof}[Proof sketch]
Property (a) is a standard Peter--Weyl computation. For (b), the mass
distribution is immediate from the amplitude weights;\ the multiplier
symbol follows by expanding $|h|^{2}$ with the SU(2) Clebsch--Gordan
rule $\chi_{j_1}\chi_{j_2} = \sum_{J = |j_1 - j_2|}^{j_1+j_2}\chi_{J}$
(unit steps in $J$, so $j_1 + j_2 + J \in \mathbb{Z}$)
and reading off the convolution eigenvalue
$c_{j'}/((2j'+1)Z)$ on each block;\ unitality at $j' = 0$ is
$c_{0} = \sum_{j}(2j+1) = Z$, and the $\nmax = 2$ values are verified
bit-exactly against independent Weyl-integration quadrature
(supporting computation in the project repository,
\texttt{debug/p45\_adversarial\_L2\_memo.md}). Property (c):\ a unital positive map on a
commutative domain is completely positive, and a UCP map has cb-norm
exactly $1$;\ on the truncated side the induced central multiplier
acts block-scalar on $\bigoplus_{\pi}\Mat_{d_{\pi}}(\C)$ with cb-norm
$\sup_{\pi}|\sigma(\pi)| = 1$ — an elementary statement requiring no
amenability input.

For (d), the closed form (d.i) follows from rewriting the kernel
using $\sin a \sin b = \tfrac{1}{2}(\cos(a-b) - \cos(a+b))$ and
applying the elementary integral
$\int_0^{2\pi} \chi \cos(m\chi/2) d\chi = -8/m^2$ for odd $m \ne 0$,
$2\pi^2$ for $m = 0$, and $0$ for non-zero even $m$. Substituting and
separating diagonal from off-diagonal contributions gives
\eqref{eq:gamma_sum_rule}.

For (d.ii), expand
$\sqrt{k_1 k_2} = \sqrt{a(a+d)} \approx a + d/2 + O(d^2/a)$ for
$k_2 = k_1 + d$, $d$ odd, $a \gg d$. The leading-order sum is
$\sum_{d\,\mathrm{odd}} 1/d^2 \cdot \sum_a a \to (\pi^2/8) \cdot
n^2/2$, giving $T_n / Z_n \to \pi^2/4$ and $\gamma_n \to 0$. The
sub-leading $\log n / n$ correction is the net of two logarithmic
pieces --- the triangle truncation of the odd-$d$ sum (coefficient
$2$) and the $d/2$ term of $\sqrt{a(a+d)}$ (coefficient $-1$) ---
computed in Appendix~\ref{app:stein_weiss}, Step~3:
\begin{equation*}
   \frac{T_n}{Z_n}
   \;=\; \frac{\pi^2}{4} \;-\; \frac{\log n}{n} \;+\; O(1/n),
\end{equation*}
which gives $\gamma_n = (4/\pi) \log n / n + O(1/n)$.

For (d.iii), \textbf{[PANEL-VERIFIED]} the function
$n \mapsto n \gamma_n / \log n$ is
monotonically decreasing for $n \ge 3$ (verified on the sample $n \in
\{3, 4, 5, 7, 10, 15, 20, 30, 50, 100, 200, 500, 1000\}$ by the frozen
test \texttt{tests/test\_central\_fejer\_su2.py};\ no proof of
monotonicity for all $n$ is given),
with maximum at $n = 2$ equal to
$2 \gamma_2 / \log 2 = 2(\pi - 64\sqrt{2}/(27\pi)) / \log 2
\approx 5.986$. The uniform bound follows.
\end{proof}

\begin{remark}[Connection to the circle Fej\'er estimate]
\label{rem:circle_fejer}
The asymptotic rate $\gamma_n \sim (4/\pi) \log n / n$ is the
non-abelian analog of the classical Fej\'er-on-the-circle estimate
\begin{equation}\label{eq:circle_fejer_moment}
   \frac{1}{2\pi}\int_{\T^1} F_n(\theta)\, |\theta|\, d\theta
   \;\sim\; \frac{2}{\pi}\,\frac{\log n}{n},
   \qquad
   F_n(\theta) = \frac{1}{n}\,\frac{\sin^2(n\theta/2)}{\sin^2(\theta/2)},
\end{equation}
for the probability-normalised Fej\'er kernel
($\frac{1}{2\pi}\int_{\T^1} F_n = 1$; Zygmund~\cite{zygmund2002}
Vol.~I, Ch.~III). The circle moment has the exact closed form
$\frac{1}{2\pi}\int_{\T^1} F_n|\theta| = \frac{\pi}{2} -
\frac{4}{\pi}\sum_{k\ \mathrm{odd},\,k<n}(1-k/n)/k^2$ (from
$\int_0^\pi\theta\cos k\theta\,d\theta = ((-1)^k-1)/k^2$), whose
$\log n/n$ coefficient is $2/\pi$ (the $\tfrac{1}{n}\sum_{k\ \mathrm{odd}}1/k
\sim \tfrac{1}{2n}\log n$ term). \textbf{[SYMBOLIC + MEASURED]} The circle
constant is $2/\pi$, not $4/\pi$:\ the doubling estimator
$(2n\,m_{2n} - n\,m_n)/\log 2$ on the closed form gives $0.63662$ at
$n = 400$ ($2/\pi = 0.63662$), the $\SU(2)$ doubling estimator of
Lemma~\ref{lem:L2}(d.ii), $a_n = (2n\gamma_{2n} - n\gamma_n)/\log 2$,
reads $1.2785$ at $n = 800$ and $1.2761$ at $n = 1600$ (i.e.\ using
$\gamma_{1600}$ and $\gamma_{3200}$; three digits of $4/\pi = 1.2732$,
approaching from above with the residual shrinking by ${\sim}1.85\times$
per doubling; Appendix~\ref{app:stein_weiss}), and the ratio of the two
estimators is $2.008$ at $n = 800$
(\texttt{tests/test\_trunk\_qa\_fejer\_4\_over\_pi.py}:\
\texttt{test\_circle\_fejer\_closed\_form\_matches\_quadrature},
\texttt{test\_circle\_fejer\_constant\_is\_2\_over\_pi},
\texttt{test\_su2\_constant\_is\_twice\_circle}). An earlier version of
this remark stated the circle constant as $4/\pi$, ``the same on both
sides''; that statement is withdrawn (2026-09-02 certification run).
\textbf{[OBSERVATION]} In the dual-Coxeter normalisation the $\SU(2)$
constant is twice the circle constant; in the unit-$\sthree$
normalisation (geodesic distance, the circle's own convention) the two
coincide at $2/\pi$ (circle:\ $2\Vol(S^0)/\Vol(S^1) = 4/(2\pi) = 2/\pi$;
$\SU(2)$:\ $\Vol(S^2)/\Vol(\SU(2)) = 4\pi/(2\pi^2) = 2/\pi$ --- the same
number, reached with a factor $2$ on the circle side and without it on
the $\SU(2)$ side, so no common ``$2\Vol(\text{base})/\Vol(\text{group})$''
form is claimed), the
ratio of the unit-metric $\SU(2)$ estimator to the circle's being
$1.004$ at $n = 800$ (2026-09-03 correction:\ the factor $2$ reported
before is the metric scale, not a group-theoretic signature). No
derivation of the $\SU(2)$ constant from the circle constant is given
here; the $\sin^2(\chi/2)$ conjugacy-class weight and the
$\sqrt{2j+1}$ Plancherel weight of Definition~\ref{def:central_fejer}
together change both the kernel and the constant, while preserving
the leading $\log n / n$ behaviour that the Ces\`aro-1 (Fej\'er)
first moment already has on the circle. The one-log slowdown quoted
in the abstract is relative to the Leimbach--van~Suijlekom torus
rate $O(1/\nmax)$~\cite{leimbach_vs2024}, not relative to the circle
Fej\'er first moment.
\end{remark}

\begin{remark}[The volume ratio behind the constant]
\textbf{[OBSERVATION]} In the dual-Coxeter normalisation the asymptotic
constant coincides numerically with
\begin{equation*}
   \frac{4}{\pi}
   \;=\; \frac{\Vol(S^2)}{\pi^2}
   \;=\; \frac{4\pi}{\pi^2}.
\end{equation*}
\emph{Which volume ratio (2026-09-03).}  The constant is a genuine quotient of sphere volumes --- but not the one its usual name suggests.  For unit-radius spheres $\Vol(S^2)/\Vol(\sthree) = 4\pi/2\pi^2 = 2/\pi$ exactly, and the dual-Coxeter constant is twice that.  It is \emph{not} the Hopf fibration's base-to-total ratio:\ the Hopf map is a Riemannian submersion $\sthree(r) \to S^2(r/2)$, so that ratio is $1/(2\pi)$ on the unit sphere and $1/(4\pi)$ in the dual-Coxeter metric.  The \emph{volume content} of the constant --- what places it in the $k = 0$ (M1) slot of the master Mellin engine and in $\mathbb{Q}[\pi, \pi^{-1}]$ --- is therefore intact;\ only the ``Hopf base'' label is a misnomer, that base carrying half the radius of the sphere whose volume appears.
Read as the unit-sphere quotient $\Vol(S^2)/\Vol(\sthree)$ rather than as a
Hopf base-to-total ratio, the constant still carries pure volume content, and is one of
the natural transcendentals on $\SU(2)$ that survives at the
leading-order rate of any Cesaro-2-style averaging procedure.
\end{remark}

% ---------------------------------------------------------------------
\subsection{Lemma L3:\ Lipschitz comparison with $C_3 = 1$}
\label{sec:L3}
% ---------------------------------------------------------------------

The next lemma is the Lipschitz comparison estimate that controls the
operator norm of the commutator $[\DCH, \Mat_f]$ in terms of the
classical Lipschitz norm of $f$ on the round $\sthree$.

\begin{lemma}[L3:\ Lipschitz comparison, truthful $\DCH$ on $\sthree$;\
\textbf{[PANEL-VERIFIED]}]
\label{lem:L3}
For every $f \in C^\infty(\sthree)$ and every $\nmax \ge 1$,
\begin{equation}\label{eq:L3_main}
   \opnorm{[\DCH,\, \Mat_f]}
   \;\le\; C_3 \cdot \norm{\nabla f}_{L^\infty(\sthree)}, \qquad
   C_3 = 1,
\end{equation}
where the round-$\sthree$ Lipschitz norm
$\norm{\nabla f}_{L^\infty}^2 = (\partial_\chi f)^2 + \sin^{-2}\chi\,
(\partial_\theta f)^2 + \sin^{-2}\chi \sin^{-2}\theta\, (\partial_\phi f)^2$
is taken in the standard polar parametrisation.
\end{lemma}

\noindent Frozen falsifier:\
\texttt{tests/test\_r25\_l3\_lipschitz\_bound.py}
(\texttt{test\_eq\_l3\_main\_inequality},
\texttt{test\_l3\_C3\_uniform\_in\_n\_max}).

\begin{proof}
The truthful $\DCH$ is diagonal in the basis
$\{\psi_{n, l, m_j, \chi}\}$ with eigenvalue $\chi(n + 1/2)$, so
\begin{equation}\label{eq:commutator_shell_diff}
   \bigl[\DCH, \Mat_f\bigr]_{(n', l', m'_j, \chi'),\,(n, l, m_j, \chi)}
   \;=\; \delta_{\chi, \chi'}\,\chi\,(n - n')\,(\Mat_f^{\mathrm{Weyl}})
   _{(n', l', m'_j),\,(n, l, m_j)}.
\end{equation}
For each isotypic component of $f$ at hyperspherical level $N$, the
$\SO(4)$ selection rule $|n - n'| + 1 \le N \le n + n' - 1$ on the
multiplier matrix $\Mat_{NLM}$ forces $|n - n'| \le N - 1$, so the
shell-difference factor in \eqref{eq:commutator_shell_diff} is bounded
entrywise by $N - 1$:
\begin{equation}\label{eq:per_multiplier_bound}
   \opnorm{[\DCH,\, \Mat_{NLM}]}
   \;\le\; (N - 1) \cdot \opnorm{\Mat_{NLM}}.
\end{equation}

The Laplace--Beltrami spectrum on $\sthree$ at level $N$ is $-(N^2 -1)$,
so the round-$\sthree$ Lipschitz norm of a unit-normalised harmonic
$Y^{(3)}_{NLM}$ scales as
\begin{equation}\label{eq:lip_scaling}
   \norm{\nabla Y^{(3)}_{NLM}}_{L^\infty}
   \;\sim\; \sqrt{N^2 - 1}\;\cdot\;\Vert \Mat_{NLM}\Vert_{\mathrm{rel}},
\end{equation}
where \textbf{[PANEL-VERIFIED]} the relative-multiplier-norm prefactor
is bounded above by $1$
on the unit-normalised panel\footnote{This is verified numerically on
the test panel $\{Y^{(3)}_{NLM} : 2 \le N \le 4\}$ in supporting
computations only; a general proof of the inequality
$\norm{\nabla Y}_{L^\infty} \le \sqrt{N^2-1}\,\Vert \Mat_{NLM}\Vert_{\mathrm{rel}}$
is open (the natural Sobolev route needs $H^s$ with $s > 3/2$ on
$\sthree$, not $H^1$, so the low-harmonic smoothness must enter
quantitatively).  Because $C_3 \le 1$ is derived through
\eqref{eq:lip_scaling}, the $C_3 \le 1$ statement is
panel-verified on the stated test set rather than proved in general;
the unconditional main theorem does not use Lemma L3 and is
unaffected.}.

Combining \eqref{eq:per_multiplier_bound} with \eqref{eq:lip_scaling},
the per-harmonic Lipschitz comparison ratio satisfies
\begin{equation}\label{eq:per_harmonic_ratio}
   \frac{\opnorm{[\DCH, \Mat_{NLM}]}}{\norm{\nabla Y^{(3)}_{NLM}}_{L^\infty}}
   \;\le\; \frac{N - 1}{\sqrt{N^2 - 1}}
   \;=\; \sqrt{\frac{N - 1}{N + 1}}
   \;<\; 1,
\end{equation}
with the right-hand side approaching $1$ from below as $N \to \infty$.
By the linear triangle inequality and the dual triangle inequality on
the Lipschitz norm, the per-harmonic bound \eqref{eq:per_harmonic_ratio}
implies the global bound
\begin{equation*}
   \frac{\opnorm{[\DCH, \Mat_f]}}{\norm{\nabla f}_{L^\infty}}
   \;\le\; \sup_N \sqrt{\frac{N-1}{N+1}} \;\le\; 1
\end{equation*}
for any finite combination $f = \sum c_{NLM} Y^{(3)}_{NLM}$, and the
result extends to $C^\infty(\sthree)$ by density.
\end{proof}

\begin{remark}
\textbf{[PANEL-VERIFIED]} The bound $C_3 = 1$ is approached as
$N \to \infty$ and is strictly below $1$ at every finite cutoff
$\nmax$;\ we do not claim it is attained.  Specifically,
$C_3(\nmax) \le \sup_{N \le \nmax} \sqrt{(N-1)/(N+1)}
= \sqrt{(\nmax - 1)/(\nmax + 1)}$ on the sub-envelope $N \le \nmax$
--- the multiplier substrate in fact reaches $N \le 2\nmax - 1$
(Definition~\ref{def:berezin}), so the per-cutoff values listed next
bound the sub-envelope only.  That value is $0$ at $\nmax = 1$,
$1/\sqrt{3} \approx 0.577$ at $\nmax = 2$, $1/\sqrt{2} \approx 0.707$
at $\nmax = 3$, $\sqrt{3/5} \approx 0.775$ at $\nmax = 4$, and so on.
Taking the supremum over the full substrate range instead would give
$\sqrt{1 - 1/\nmax}$.  \emph{That expression must not be substituted
here without a measurement.}  It is textually identical to a constant
retired in 2026-06 (Papers~45/46): the \emph{operator-norm}-normalised
envelope $\opnorm{[D_{\mathrm{GV}}, \Mat]}/\opnorm{\Mat} \le
\sqrt{1-1/\nmax}$, which is false on the operator system.  The bound
\eqref{eq:per_harmonic_ratio} above is normalised by
$\norm{\nabla Y}_{L^\infty}$, not by $\opnorm{\Mat}$, and is a
different statement that survives that retraction.  The correct
gradient-normalised per-cutoff envelope is an open measurement, not an
inference from the substrate range;\ the load-bearing conclusion
$C_3 \le 1$ is unaffected either way.
\end{remark}

% ---------------------------------------------------------------------
\subsection{Lemma L4:\ Berezin reconstruction}
\label{sec:L4}
% ---------------------------------------------------------------------

\begin{definition}[Berezin map, convolution form]\label{def:berezin}
For $\nmax \ge 1$ and $f \in C^\infty(\sthree)$, define
\begin{equation}\label{eq:B_def}
   B_{\nmax}(f)
   \;:=\; P_{\nmax}\,\Mat_{\KS_{\nmax} \ast f}\,P_{\nmax}
   \;=\; \sum_{N \le 2\nmax - 1}\sum_{L=0}^{N-1}\sum_{|M| \le L}
   \sigma_{\nmax}(N)\, c_{NLM}\, \Mat_{NLM},
\end{equation}
where $\sigma_{\nmax}(N)$ is the multiplier symbol of
Lemma~\ref{lem:L2}(b) under the bijection $n = 2j+1$ and
$f = \sum c_{NLM} Y^{(3)}_{NLM}$. (A sum form weighted by the mass distribution $N/Z_{\nmax}$ in
place of $\sigma_{\nmax}(N)$ would not be unital —
$B(\mathbf{1}) = Z^{-1} I$ — and would fail the approximate-identity
property maximally at $f = \mathbf{1}$. Note the image reaches
multiplier labels up to the achievable envelope $N \le 2\nmax - 1$.)
\end{definition}

\begin{lemma}[L4:\ properties of the Berezin reconstruction]\label{lem:L4}
The map $B_{\nmax}: C^\infty(\sthree) \to \Op_{\nmax}$ has the
following four properties:
\begin{enumerate}\setlength{\itemsep}{2pt}
\item[(a)] (Positivity.) If $f \ge 0$ on $\sthree$, then $B_{\nmax}(f)$
is positive semi-definite.
\item[(b)] (Contractivity.) $\opnorm{B_{\nmax}(f)} \le \norm{f}_{L^\infty}$.
\item[(c)] (Approximate identity.)
$\opnorm{B_{\nmax}(f) - P_{\nmax} \Mat_f P_{\nmax}}
\le \gamma_{\nmax} \cdot \norm{\nabla f}_{L^\infty}$.
\item[(d)] (Compatibility with L3.) \textbf{[PANEL-VERIFIED]}
$\opnorm{[\DCH, B_{\nmax}(f)]} \le \norm{\nabla f}_{L^\infty}$
(inherits the panel scope of Lemma~\ref{lem:L3} through $C_3 = 1$).
\end{enumerate}
\end{lemma}

\begin{proof}
\textbf{(a)} Equivalent to the convolution form
$B_{\nmax}(f) = P_{\nmax} (\KS_{\nmax} * f) P_{\nmax}$:\ since
$\KS_{\nmax} \ge 0$ pointwise (Lemma~\ref{lem:L2}(a)) and the
convolution of two non-negative functions on a compact group is
non-negative, $\KS_{\nmax} * f \ge 0$. Compression by an orthogonal
projection preserves positivity:\ for any $\psi \in \Hilb_{\nmax}$,
\begin{equation*}
\langle \psi | P_{\nmax} \Mat_g P_{\nmax} | \psi \rangle
\;=\; \int_{\sthree} g(\omega)\,|P_{\nmax} \psi(\omega)|^2\,d\Omega \;\ge\; 0.
\end{equation*}

\textbf{(b)} By Young's inequality,
$\norm{\KS_{\nmax} * f}_{L^\infty} \le \norm{\KS_{\nmax}}_{L^1}
\norm{f}_{L^\infty} = \norm{f}_{L^\infty}$. Compression contracts
operator norms, so $\opnorm{B_{\nmax}(f)} \le \norm{f}_{L^\infty}$.

\textbf{(c)} The standard Lipschitz/good-kernel estimate gives
\begin{equation*}
   \norm{\KS_{\nmax} * f - f}_{L^\infty}
   \;\le\; \int_{\SU(2)} \KS_{\nmax}(g)\,|f(g \cdot x) - f(x)|\,dg
   \;\le\; \gamma_{\nmax}\,\norm{\nabla f}_{L^\infty},
\end{equation*}
which is the SU(2) version of the standard approximate-identity
estimate (Stein--Weiss~\cite{stein_weiss1971}, Ch.~I).
Compression then gives
$\opnorm{B_{\nmax}(f) - P_{\nmax} \Mat_f P_{\nmax}}
\le \gamma_{\nmax}\norm{\nabla f}_{L^\infty}$.

\textbf{(d)} The cleanest argument goes via the convolution form
$B_{\nmax}(f) = P_{\nmax}\,\Mat_{\KS_{\nmax} * f}\,P_{\nmax}$ used in
parts (a) and (c). Let $g := \KS_{\nmax} * f \in C^\infty(\sthree)$.
Since $\KS_{\nmax} \ge 0$ and
$\norm{\KS_{\nmax}}_{L^1} = 1$ by Lemma~\ref{lem:L2}(a), Young's
inequality for the gradient
\begin{equation}\label{eq:young_gradient}
   \norm{\nabla g}_{L^\infty}
   \;=\; \norm{\nabla(\KS_{\nmax} * f)}_{L^\infty}
   \;=\; \norm{\KS_{\nmax} * \nabla f}_{L^\infty}
   \;\le\; \norm{\KS_{\nmax}}_{L^1}\,\norm{\nabla f}_{L^\infty}
   \;=\; \norm{\nabla f}_{L^\infty}
\end{equation}
gives a smoothed function $g$ whose Lipschitz norm is no larger than
that of $f$ (left-translation by $\KS_{\nmax}$ commutes with the
gradient on the bi-invariant metric and convolution with a probability
density does not increase the Lipschitz norm). Now apply
Lemma~\ref{lem:L3} to $g$ in the truncated multiplier representation:
$B_{\nmax}(f) = \Mat_g$ in $\Op_{\nmax}$, so
\begin{equation*}
   \opnorm{[\DCH, B_{\nmax}(f)]}
   \;=\; \opnorm{[\DCH, \Mat_g]}
   \;\le\; C_3\,\norm{\nabla g}_{L^\infty}
   \;\le\; C_3\,\norm{\nabla f}_{L^\infty}
   \;=\; \norm{\nabla f}_{L^\infty},
\end{equation*}
where the final equality uses $C_3 = 1$ from Lemma~\ref{lem:L3}.
\end{proof}

\begin{remark}[Non-K\"ahler analog]
The Berezin reconstruction \eqref{eq:B_def} differs from the standard
Hawkins~\cite{hawkins2000} construction on K\"ahler manifolds, which
proceeds via the holomorphic Toeplitz form on a Hardy space. Since
$\sthree = \SU(2)$ admits no compatible almost-complex structure
(it is parallelisable but not K\"ahler), the natural analog is the
central-spectral-Fej\'er-kernel form, which is the
abelian-compact-group version
of the same idea applied to the conjugacy-class measure on $\SU(2)$.
The K\"ahler analog on $S^2 = \SU(2)/U(1)$ via Berezin--Toeplitz on
the Hopf base has been worked out independently
(see~\cite{ucp_maps_2024} and references therein) and is consistent
with the present construction in the appropriate fibration limit.
\end{remark}

% ---------------------------------------------------------------------
\emph{Panel note (2026-09-03).}  The reference implementation
(\texttt{geovac/gh\_convergence.py}) reports reach and height on a fixed
test panel of low-degree harmonics.  Those panel numbers do \emph{not}
verify this lemma:\ the Lipschitz-normalised panel height rises with the
cutoff while $\gamma_{\nmax}$ falls, and the two cross between
$\nmax = 6$ and $\nmax = 7$ (measured 2026-09-03):
\[
  \begin{array}{c|cccccc}
    \nmax & 2 & 3 & 4 & 5 & 6 & 7 \\ \hline
    \gamma_{\nmax} & 2.075 & 1.610 & 1.322 & 1.130 & 0.990 & 0.883 \\
    \text{height}/\|f\|_{\mathrm{Lip}} & 0.667 & 0.808 & 0.878 & 0.913 & 0.937 & 0.952 \\
    \text{margin} & +1.408 & +0.803 & +0.445 & +0.217 & +0.053 & -0.069
  \end{array}
\]
The panel is not the unit-Lipschitz ball quantified over here.  The proof
below does not use the panel;\ identifying the panel-side quantity is an
open check.

\subsection{Lemma L5:\ state-space Gromov--Hausdorff assembly}
\label{sec:L5}
% ---------------------------------------------------------------------

A remark on framework before assembling. In Latr\'emoli\`ere's
quantum Gromov--Hausdorff propinquity for metric spectral
triples~\cite{latremoliere2016, latremoliere_metric_st_2017}, the
distance is an infimum over tunnels carried by quantum metric spaces
with quantum isometries, under hypotheses (Leibniz Lip-norms on
C$^{*}$-algebras;\ kernel condition) not established for the
truncated operator system;\ a pair of UCP maps is not a tunnel there.
The distance the pair $(B_{\nmax}, P_{\nmax})$ certifies is
van~Suijlekom's state-space Gromov--Hausdorff
distance~\cite{vs2021_jgp}, via two-directional approximation data in
the Leimbach--van~Suijlekom pattern~\cite{leimbach_vs2024}. We write
\begin{equation}\label{eq:propinquity_def}
   \Lambda(\Tcal_1, \Tcal_2)
   \;:=\; \mathrm{d}_{\mathrm{GH}}^{\,\mathrm{vS}}\!\bigl(
      S(\Op_{1}),\, S(\Op_{2})\bigr)
\end{equation}
for this distance (between state spaces with their respective
Monge--Kantorovich metrics), and retain the constituent names
reach$_{B}$, reach$_{P}$, height$_{B}$, height$_{P}$ for the four
approximation estimates the pair supplies. The truncated-side
Monge--Kantorovich metric is non-degenerate unconditionally under the
translation-seminorm metrization
(Proposition~\ref{prop:kernel_condition});\ the distance bound itself
is Theorem~\ref{thm:main_unconditional}, and the pair below supplies
the explicit forward-reconstruction refinement.

\begin{lemma}[L5:\ Berezin-pair refinement of the forward
direction;\ \textbf{[PANEL-VERIFIED]}]\label{lem:L5}
Within the framework of \S~\ref{sec:named_gaps} (kernel condition
Proposition~\ref{prop:kernel_condition};\ distance bound
Theorem~\ref{thm:main_unconditional}), the pair
$(B_{\nmax}, P_{\nmax})$ with $B_{\nmax}$ from
Definition~\ref{def:berezin} and $P_{\nmax}$ the standard spectral
projection yields the bound
\begin{equation}\label{eq:L5_bound}
   \Lambda(\Tcal_{\nmax}, \Tcal_{\sthree})
   \;\le\; C_3 \cdot \gamma_{\nmax},
\end{equation}
with $C_3 = 1$ (Lemma~\ref{lem:L3}) and $\gamma_{\nmax}$ from
Lemma~\ref{lem:L2}(d).
\end{lemma}

\begin{proof}
The four constituents a Latr\'emoli\`ere tunnel would require are:
\begin{enumerate}\setlength{\itemsep}{2pt}
\item $\mathrm{reach}_B
:= \sup_{\Lipnorm{f} \le 1}
\opnorm{B_{\nmax}(f) - P_{\nmax} \Mat_f P_{\nmax}}$,
\item $\mathrm{reach}_P$, the dual reach via the partial inverse,
\item $\mathrm{height}_B$, the Lipschitz-distortion envelope of
$B_{\nmax}$ (defined precisely below),
\item $\mathrm{height}_P = 0$ for the orthogonal projection.
\end{enumerate}

\emph{Reach.} By L4(c), $\mathrm{reach}_B \le
\gamma_{\nmax}\cdot \norm{\nabla f}_{L^\infty} \le
\gamma_{\nmax}$ on the unit Lipschitz ball
$\{\Lipnorm{f} \le 1\}$, where \textbf{[PANEL-VERIFIED]} we have used
L3 to convert
$\norm{\nabla f}_{L^\infty} \le \Lipnorm{f}$ with constant $C_3 = 1$,
so this route --- unlike Theorem~\ref{thm:main_unconditional} ---
carries Lemma~\ref{lem:L3}'s panel scope.
\textbf{(Dissolved.)} The dual estimate is no longer needed:\
Theorem~\ref{thm:main_unconditional} obtains the distance bound via
the lifted-state map, whose almost-inverse defect is the conjugation
average $\Phi$ at the same moment $\gamma_{\nmax}$
(Remark~\ref{rem:no_inverse}) — no partial inverse of $B_{\nmax}$,
and no transference estimate, enters anywhere. (For the record:\ a
partial inverse of $B_{\nmax}$ on the full support window would be
governed by the symbol \emph{minimum} and grow like $O(\nmax^2)$,
which is why no forward bound could supply it.) The pair $(B_{\nmax},
P_{\nmax})$ retains its role as the explicit forward-reconstruction
refinement.

\emph{Height.} The height constituent is modeled on the propinquity
for metric spectral triples \cite{latremoliere_metric_st_2017}, which
measures the Lipschitz-distortion of the tunnel maps with respect to the
Lipschitz seminorms, \emph{not} the operator norm of the maps
themselves. Specifically,
\begin{equation}\label{eq:height_B_definition}
   \mathrm{height}_B
   \;:=\; \sup_{\Lipnorm{f} \le 1}
   \bigl|\,\Lipnorm{f} - \Lipnorm{B_{\nmax}(f)}^{\Op_{\nmax}}\bigr|.
\end{equation}
By L4(d), $\Lipnorm{B_{\nmax}(f)}^{\Op_{\nmax}} =
\opnorm{[\DCH, B_{\nmax}(f)]} \le \norm{\nabla f}_{L^\infty}
\le \Lipnorm{f}$, so the absolute difference in
\eqref{eq:height_B_definition} is bounded above by
\begin{equation*}
   \Lipnorm{f} - \Lipnorm{B_{\nmax}(f)}^{\Op_{\nmax}}
   \;=\; \Lipnorm{f} - \opnorm{[\DCH, \Mat_{\KS_{\nmax}*f}]}
   \;\le\; \Lipnorm{f} - \Lipnorm{\KS_{\nmax}*f}.
\end{equation*}
By the standard Lipschitz/good-kernel estimate (the same
approximate-identity inequality used in L4(c)), the difference of
Lipschitz norms is
bounded by $\gamma_{\nmax}$:\ $\Lipnorm{f} - \Lipnorm{\KS_{\nmax}*f}
\le \gamma_{\nmax}\Lipnorm{f}$. Hence
\begin{equation*}
   \mathrm{height}_B \;\le\; \gamma_{\nmax}, \qquad
   \mathrm{height}_P \;=\; 0.
\end{equation*}

Inserting into \eqref{eq:propinquity_def}:
\begin{equation*}
   \Lambda(\Tcal_{\nmax}, \Tcal_{\sthree})
   \;\le\; \max\bigl(\gamma_{\nmax}, \gamma_{\nmax},
   \gamma_{\nmax}, 0\bigr)
   \;=\; \gamma_{\nmax}
   \;=\; C_3 \cdot \gamma_{\nmax}
   \;\to\; 0 \;\;\text{as}\;\; \nmax \to \infty,
\end{equation*}
which gives \eqref{eq:L5_bound}.
\end{proof}

\begin{remark}[On the height term]\label{rem:height_definition}
The height in \eqref{eq:height_B_definition} is the
spectral-triple-level analog of the height in
Latr\'emoli\`ere~\cite[\S 4]{latremoliere_metric_st_2017}, which
measures the Lipschitz-seminorm-comparison error introduced by the
tunnel maps. The naive bound $\opnorm{B_{\nmax}(f)} \le \norm{f}_{L^\infty}
\le \pi$ (from L4(b) plus the $\sthree$ diameter) is not the relevant
height for the spectral-triple propinquity:\ the operator norm of
$B_{\nmax}(f)$ enters the height for the quantum compact metric
space propinquity \cite{latremoliere2016} but not for the spectral
triple version taken as the model here. The reduction of the height to
$\gamma_{\nmax}$ via the Lipschitz-comparison form is what gives the
$n$-rate-controlled state-space GH bound of Theorem~\ref{thm:main}.
\end{remark}

% =====================================================================
\section{Main theorem}\label{sec:main_theorem}
% =====================================================================

\begin{theorem}[Main theorem;\ rate-controlled restatement]
\label{thm:main}
\textbf{[INTERNAL THEOREM]}
Unconditionally (Theorem~\ref{thm:main_unconditional}), the
truncated Camporesi--Higuchi spectral triples $\Tcal_{\nmax}$
converge to the round-$\sthree$ Camporesi--Higuchi spectral triple
$\Tcal_{\sthree}$ in van~Suijlekom's state-space Gromov--Hausdorff
distance~\cite{vs2021_jgp}, under the translation-seminorm
metrization of \S~\ref{sec:named_gaps}:
\begin{equation}\label{eq:main_thm}
   \Lambda(\Tcal_{\nmax}, \Tcal_{\sthree})
   \;\le\; C_3 \cdot \gamma_{\nmax}, \qquad
   \gamma_{\nmax} = \frac{4 \log \nmax}{\pi\,\nmax}
   + O\!\bigl(1/\nmax\bigr).
\end{equation}
\textbf{[OBSERVATION]} The asymptotic constant
$4/\pi = \Vol(S^2)/\pi^2$ is the
Hopf-base measure factor in the standard $\SU(2)/U(1)$ Haar
normalisation. \textbf{[PANEL-VERIFIED]} A uniform bound
$\gamma_{\nmax} \le 6 \log \nmax / \nmax$
holds for all $\nmax \ge 2$.
\end{theorem}

\begin{proof}
Lemma~\ref{lem:L1prime} supplies the operator-system substrate;
Lemma~\ref{lem:L2} the central spectral Fej\'er kernel with closed-form
$\gamma_{\nmax}$ and asymptotic rate; Lemma~\ref{lem:L3} the Lipschitz
comparison constant $C_3 = 1$; Lemma~\ref{lem:L4} the Berezin
reconstruction with the four properties a Latr\'emoli\`ere
tunneling pair would require; and Lemma~\ref{lem:L5} the assembly into the
state-space GH bound. Combining the bound \eqref{eq:L5_bound} with the
asymptotic and uniform rates of Lemma~\ref{lem:L2}(d.ii)--(d.iii)
gives \eqref{eq:main_thm}.  This lemma-assembly route inherits the
panel scope of Lemma~\ref{lem:L3};\ the unconditional content of the
statement is Theorem~\ref{thm:main_unconditional}, which does not use
Lemma~\ref{lem:L3} at all.
\end{proof}

A companion proposition identifies the state-space GH limit explicitly.

\begin{proposition}[Limit identification]
\label{prop:limit_identification}
\textbf{[INTERNAL THEOREM]} The state-space GH limit of the sequence
$(\Tcal_{\nmax})_{\nmax \ge 1}$ is
\begin{equation*}
   \lim_{\nmax \to \infty} S(\Op_{\nmax})
   \;=\; \bigl(P(\sthree),\; d_{\mathrm{Wass}}\bigr),
\end{equation*}
where $P(\sthree)$ is the space of Borel probability measures on
$\sthree$ and $d_{\mathrm{Wass}}$ is the Wasserstein--Kantorovich
distance with respect to the round-$\sthree$ geodesic distance.
\end{proposition}

\begin{proof}[Proof sketch]
Theorem~\ref{thm:main} guarantees $\Lambda(\Tcal_{\nmax},
\Tcal_{\sthree}) \to 0$ in the state-space GH distance, so the limit
is the continuum state space. By Kantorovich--Rubinstein duality, the
Connes distance on $\Tcal_{\sthree}$ coincides with
$d_{\mathrm{Wass}}$ on $P(\sthree)$. The L4(c) approximate-identity
property pins this to the actual limit (rather than an isometry
class). This is the same limit identification as in
van~Suijlekom~\cite{vs2021_jgp} and
Gaudillot-Estrada--van~Suijlekom~\cite{gaudillot_vs2025} for the
scalar case.
\end{proof}

% =====================================================================
\section{Discussion}\label{sec:discussion}
% =====================================================================

\subsection{Why $\SU(2)$ was harder than the torus}\label{sec:why_su2}

The flat-torus result of Leimbach--van~Suijlekom~\cite{leimbach_vs2024}
proceeds by Schur multiplier analysis on the abelian dual $\Z^d$,
where the Fej\'er kernel is a tensor product of one-dimensional
Fej\'er kernels and the cb-norm is computed via the standard
abelian-group identity $\Vert T\Vert_{\mathrm{cb}} = \norm{\widehat{T}}_\infty$.
Three substantive complications arise on $\SU(2)$:

\textbf{(i) Non-abelian Peter--Weyl.} On $\SU(2)$ the Plancherel
decomposition is not over $\Z^d$ but over half-integer weights
$j \in \tfrac{1}{2}\N_{\ge 0}$ with isotypic blocks of dimension
$2j+1$. The natural Plancherel weight $\sqrt{2j+1}$ is required to
match the Fock-shell Peter--Weyl bijection $n = 2j+1$, and this is
what produces the natural-coefficient kernel of
Definition~\ref{def:central_fejer}. (For central multipliers the
cb-norm computation is elementary — block-scalar action — on any
compact group;\ the genuine non-abelian difficulty is the
Clebsch--Gordan structure of the multiplier symbol, not the cb-norm.)

\textbf{(ii) No K\"ahler structure on $\sthree$.} The standard Berezin
reconstruction of Hawkins~\cite{hawkins2000} uses the holomorphic
Toeplitz form on a Hardy space, which requires a compatible
almost-complex structure. $\sthree = \SU(2)$ has no such structure
(this is the standard fact that $\sthree$ has trivial tangent bundle
but no integrable almost-complex structure). The natural alternative
is the central-spectral-Fej\'er-kernel form
\eqref{eq:B_def}, which is the abelian-compact-group analog
of the Connes--van~Suijlekom~\cite{connes_vs2021} Fej\'er
reconstruction on $S^1$, lifted to the non-abelian setting via the
Peter--Weyl decomposition.

\textbf{(iii) The $\SO(4)$ selection rule.} \textbf{[OBSERVATION]} The
Lipschitz comparison
\eqref{eq:per_multiplier_bound} is as tight as it is because of the
multiplier
sparsity enforced by the $\SO(4)$ selection rule on the Avery 3-Y
integral;\ we do not prove the bound is attained. The torus analog has no comparable sparsity:\ on $\T^d$,
multiplier matrices are full Toeplitz matrices, and the Lipschitz
comparison ratio is bounded by the
top mode index $|k|$ rather than by the cleaner
$\sqrt{(N-1)/(N+1)}$ form.

The cumulative effect is that the SU(2) rate is one log factor slower
than the torus rate ($\log n / n$ vs $1 / n$). This is the standard
slowdown of non-abelian compact groups relative to flat tori in
Cesaro/Fej\'er theory. \textbf{[OBSERVATION]} In the dual-Coxeter normalisation the
asymptotic constant $4/\pi$ is twice the circle Fej\'er first-moment
constant $2/\pi$; on the unit $\sthree$ the two coincide, the factor
$2$ being the metric scale (Remark~\ref{rem:circle_fejer},
2026-09-03); we read the $\SU(2)$ value as the Hopf-base measure
factor $\Vol(S^2)/\Vol(\SU(2))$ (unit $\sthree$), i.e.\
$2\Vol(S^2)/\Vol(\SU(2))$ in the dual-Coxeter normalisation. (An earlier version of this
sentence stated the constant as ``identical on both sides''; withdrawn
2026-09-02.)

\subsection{Marcolli--van~Suijlekom gauge-network framework}
\label{sec:marcolli_vs}

The convergence theorem
\ref{thm:main} provides the metric foundation for almost-commutative
spectral triple constructions hosted on $\sthree$. The
Marcolli--van~Suijlekom gauge-network
framework~\cite{marcolli_vs2014}, with the
Perez-Sanchez~\cite{perez_sanchez2025} continuum-limit
correction, builds finite spectral triples from algebra-valued data on
graphs;
\textbf{[INTERNAL THEOREM]} the present convergence result identifies
one such construction --- the
Connes--van~Suijlekom truncation of $\Tcal_{\sthree}$ --- as
metric-asymptotically the round-$\sthree$ Camporesi--Higuchi
geometry. This anchors the gauge-network framework to a continuum
geometry whose Dirac spectrum is independently known and whose
spectral-action expansion is computable.

The natural next question --- whether the almost-commutative extension
$\Tcal_{\sthree} \otimes \Tcal_F$ for a finite spectral triple
$\Tcal_F$ admits an analogous convergence statement --- is open.
\textbf{[MEASURED]} Inner-fluctuation analysis of the truncated AC
extension at finite
$\nmax$ has been carried out in supporting work by the present
author and reproduces the expected gauge content under the
Marcolli--van~Suijlekom mechanism (\texttt{tests/test\_almost\_commutative.py},
Sprint H1; Paper~32 \S\,Sprint H1), but the present paper restricts
attention to the geometric (single-factor) state-space GH statement.

\subsection{Obstruction to a Riemannian--Hardy tensor product}
\label{sec:obstruction}

A natural extension of Theorem~\ref{thm:main} would aim at the
tensor-product spectral triple $\Tcal_{\sthree} \otimes
\Tcal_{S^5}^{\mathrm{Bargmann}}$ obtained by tensoring the present
Camporesi--Higuchi triple on $\sthree$ with a Bargmann-segal-type
spectral triple on $S^5$ in the holomorphic Hardy sector.
\textbf{[OBSERVATION]} This
extension is structurally obstructed at the framework level --- the
obstruction is argued from the categorical mismatch below, not proved
as a theorem. The
$\sthree$ factor is Riemannian with KO-dimension $3$ and a
second-order Laplace--Beltrami structure; the $S^5$ factor is a
Hardy-sector first-order complex Euler operator
construction~\cite{paper24} which does not fit standard NCG Dirac
axioms. The two factors live in different categories of NCG
construction, and there is no published prescription for tensoring
them. We have flagged this elsewhere as the four-layer
Coulomb/HO asymmetry resurfacing at the spectral-triple level. A
proper Riemannian--Hardy propinquity theory is not in the literature
and would require new framework development beyond the scope of the
present paper.

\subsection{Open questions}\label{sec:open}

\textbf{(i) Quantitative tightness.} Theorem~\ref{thm:main} gives the
asymptotic constant $4/\pi$ in $\gamma_{\nmax} \sim
(4/\pi) \log \nmax / \nmax$ but not the next-order constant.
\textbf{[MEASURED]} Numerical doubling-estimator data suggests
$\nmax \gamma_{\nmax} = (4/\pi) \log \nmax + b + O(1/\nmax)$ with
$b \approx 4.10$, whose closed form is undetermined at the level of
the present analysis.

\emph{Update (2026-05-16):} \textbf{[MEASURED]} the value $b$ (rotation-angle
normalisation) has been refined to $22$ decimal digits (no frozen test
in this repository computes $b$ beyond six digits; the figure is a
Paper~40 cross-reference) as
$b = 4.1093214674898756204229\ldots$
(the original $80$-digit string circulated in working memos prior to
this paper was contaminated from digit $12$ onward by float-precision
panel points; corrected via re-extraction at $\ge 400$~dps in a
follow-on sprint, see~\cite{paper40_unified} for details). PSLQ
identification has been attempted at coefficient ceiling $10^8$
against an extended L-derivative basis ($\zeta'(s), \beta'(s),
\eta'(s), L'(2, \chi_q)$, Stieltjes constants, Glaisher--Kinkelin
$\log A$) with NULL result. \textbf{[OPEN]} The closed form of $b$
remains open, and the constant joins a
small documented list of natural-but-unidentified constants in this
programme, alongside the numerical coincidence formula studied
in~\cite{paper2}.

\textbf{(ii) Lorentzian extension.} The Camporesi--Higuchi spectral
triple is Riemannian, with KO-dimension $3 \pmod 8$. On the
Riemannian side, the modular-Hamiltonian structure of the truncations
is closed at the operator-system level:\ \textbf{[MEASURED]} the
modular Hamiltonian on
$\mathcal{T}_{\nmax}$ satisfies the period closure
$\sigma_{2\pi}(O) = O$ bit-exactly at $\nmax \in \{2, 3, 4, 5\}$ for
the hemispheric-wedge-restricted state and all four physical
witnesses (Bisognano--Wichmann, Hartle--Hawking, Sewell, Unruh), in
both the geometric realization ($K = J_{\mathrm{polar}}$, integer
spectrum, so $e^{i \cdot 2\pi n} = 1$ exactly) and the
Tomita--Takesaki realization ($K_{\mathrm{TT}} = -\log\Delta$ on the
GNS Hilbert--Schmidt space), which are flow-conjugate,
$\sigma_t^{\mathrm{TT}} = \sigma_{-t}^{\alpha}$;\ see the four-witness
companion paper~\cite{paper42} and \cite{paper32}~\S~VIII.F. \textbf{[OPEN]} The
genuinely Lorentzian extension at signature $(3, 1)$ via Krein-space
spectral triples is \emph{open}, and is now known to be obstructed
along its most natural route:\ compressing a Krein-self-adjoint
product Dirac to the Krein-positive subspace annihilates the spatial
Dirac, so the compressed Lipschitz seminorm vanishes identically on
the natural operator system and no quantum-metric structure survives.
The companion Lorentzian paper~\cite{paper45} proves this degeneracy
theorem and states the repaired target (a Toeplitz temporal
multiplier algebra in the Connes--van~Suijlekom circle pattern, and a
Krein-compatible Lipschitz seminorm that does not compress the
Dirac).

\textbf{(iii) Higher-rank compact non-abelian cases.} The present
proof transcribes mechanically to $\SU(n)$ for $n \ge 3$ at the
algebraic level (Peter--Weyl, central characters, Plancherel weights)
but the Lipschitz comparison estimate
\eqref{eq:lip_scaling} requires a higher-rank version of the
Avery 3-Y integral and the corresponding selection rule.

\emph{Update (2026-05-16):} this open question has been resolved by
Paper~40~\cite{paper40_unified}, which generalises
Theorem~\ref{thm:main} to all compact connected semisimple Lie groups
with bi-invariant metric (a circle factor carries the constant $2/\pi$
of Remark~\ref{rem:circle_fejer}, so tori are excluded). The higher-rank analog of the Avery 3-Y selection
rule is the \emph{Dirac-triangle inequality}
$|D(\pi) - D(\pi')| \le \sqrt{C(\sigma)}$ for every
$\sigma \subset \pi \otimes \pi'^*$, which generalises Paper~38's
$|n - n'| \le N - 1$ shell-difference rule (the literal "Casimir
triangle inequality" $|C(\pi) - C(\pi')| \le C(\sigma)$ is false in
general; the Dirac eigenvalue, not the Casimir, is the right
ingredient). \textbf{[PANEL-VERIFIED]} The rate is \emph{not}
rank-dependent as anticipated
above:\ the rate constant $4/\pi$ (dual-Coxeter normalisation) is reported \emph{universal} across
the
class, by a Plancherel-weight $\times$ Vandermonde-Jacobian
cancellation theorem (Paper~40~\cite{paper40_unified} \S 3.2) --- rigorous at rank~$1$
($\SU(2)$, the present paper), with the rank-$\ge 2$ leg corroborated
on reduced panels and a fully rank-uniform analytical proof recorded
as a named gap in Paper~40. \textbf{[INTERNAL THEOREM]}
Asymptotic-tight
$C_3(G) = 1$ at all ranks follows from the PRV-summand bound (Kumar
1988, Invent.\ Math.\ 93) together with a Brauer--Klimyk signed-sum
closure for the interior summands, as proved in Paper~40 \S 3.3
(Lemma~L3 there; the attribution to a 1990 Vinberg paper carried no
locator and is withdrawn, 2026-09-03).

% =====================================================================
\appendix
\section{Detailed derivation of Lemma~\ref{lem:L2}(d.ii)}
\label{app:stein_weiss}
% =====================================================================

\textbf{[INTERNAL THEOREM + MEASURED]} This appendix supplies the
sum-rule / Euler--Maclaurin derivation (the ``Stein--Weiss
sharpening'' in the corpus's naming, after the real-line kernel
estimates of~\cite{stein_weiss1971}) of the asymptotic
$\lim_{n\to\infty} n\gamma_n / \log n = 4/\pi$ that appears as a
sketch in the proof of Lemma~\ref{lem:L2}. Step~3 isolates the two
sub-coefficients ($2$ and $-1$) whose net is the logarithm's
coefficient; each sub-coefficient and the net are pinned by
independent doubling estimators (test-backed), which is what closes
the constant. The argument is a standard asymptotic analysis of the
closed-form sum-rule \eqref{eq:gamma_sum_rule}; we record it here in
full because the referee may want to see the leading-order constant
pinned without having to reconstruct it.

\paragraph{Setup.}
Recall the closed form of Lemma~\ref{lem:L2}(d.i):
\begin{equation}\label{eq:gamma_closed_form_repeat}
   \gamma_n \;=\; \pi \;-\; \frac{4 T_n}{\pi Z_n}, \qquad
   Z_n = \frac{n(n+1)}{2}, \qquad
   T_n = \!\!\!\sum_{\substack{1\le k_1, k_2 \le n \\ k_1 + k_2\,\mathrm{odd}}}
                \sqrt{k_1 k_2}\Bigl[\frac{1}{(k_1-k_2)^2} - \frac{1}{(k_1+k_2)^2}\Bigr].
\end{equation}
The asymptotic claim of Lemma~\ref{lem:L2}(d.ii) is equivalent to
\begin{equation}\label{eq:Tn_asymptotic}
   \frac{T_n}{Z_n} \;=\; \frac{\pi^2}{4} \;-\; \frac{\log n}{n}
   \;+\; O\!\Bigl(\frac{1}{n}\Bigr),
\end{equation}
since $\gamma_n = \pi - (4/\pi)(T_n/Z_n) =
(4/\pi)(\log n / n) + O(1/n)$.

\paragraph{Step 1: Diagonal-dominant decomposition.}
Write $k_2 = k_1 + d$ with $d$ odd, $d \in \{\pm 1, \pm 3, \ldots\}$,
$k_1 \in \{1, \ldots, n\}$, $k_1 + d \in \{1, \ldots, n\}$. By the
$k_1 \leftrightarrow k_2$ symmetry of $T_n$, we may take $d > 0$ and
double:
\begin{equation*}
   T_n \;=\; 2 \sum_{\substack{d > 0 \\ d\,\mathrm{odd}}}
   \sum_{a=1}^{n-d}
   \sqrt{a(a+d)}\Bigl[\frac{1}{d^2} - \frac{1}{(2a+d)^2}\Bigr].
\end{equation*}

\paragraph{Step 2: Leading-order $\pi^2/4$.}
The term $1/d^2$ dominates. Using
$\sqrt{a(a+d)} = a\sqrt{1+d/a} = a + d/2 - d^2/(8a) + O(d^3/a^2)$ and
$\sum_{a=1}^{n-d} 1 = n - d$,
\begin{equation*}
   \sum_{a=1}^{n-d} \sqrt{a(a+d)} \cdot \frac{1}{d^2}
   \;=\; \frac{1}{d^2}\Bigl[\frac{(n-d)(n-d+1)}{2} + \frac{d(n-d)}{2}
   + O(d \log n)\Bigr]
   \;=\; \frac{n^2}{2 d^2} + O\!\Bigl(\frac{n}{d}\Bigr).
\end{equation*}
Summing over odd $d$ from $1$ to $n$ and using
$\sum_{d\,\mathrm{odd}} 1/d^2 = \pi^2 / 8$
(the Catalan-Euler-style identity $\sum_{d\ge 1\,\mathrm{odd}} 1/d^2
= (1 - 2^{-2})\zeta(2) = (3/4)(\pi^2/6) = \pi^2/8$),
\begin{equation*}
   T_n \;=\; 2 \cdot \frac{n^2}{2}\cdot \frac{\pi^2}{8} +
   \text{lower order} \;=\; \frac{\pi^2 n^2}{8} + \text{lower order}.
\end{equation*}
Dividing by $Z_n = n(n+1)/2 \sim n^2/2$ gives the leading
$T_n/Z_n \to \pi^2/4$, confirming $\gamma_n \to 0$.

\paragraph{Step 3: Sub-leading $\log n$ (two sub-coefficients, $2 - 1 = 1$).}
Write $\mathrm{defect}_n := \pi^2/4 - T_n/Z_n$, so that
\eqref{eq:Tn_asymptotic} is the statement $\mathrm{defect}_n = \log n/n
+ O(1/n)$, i.e.\ that the $\log n / n$ coefficient of the defect is
exactly $1$ (and $\gamma_n = (4/\pi)\,\mathrm{defect}_n$ identically by
\eqref{eq:gamma_closed_form_repeat}). Reorganise the Step-1 double
sum by $a$, with $d$ odd and $d \le n - a$, and expand
$\sqrt{a(a+d)} = a + d/2 + r(a,d)$. Two pieces carry a logarithm:
\begin{itemize}
\item[(A)] \emph{Triangle truncation of the $d$-sum.} The
$a\cdot(1/d^2)$ piece is $2\sum_{a=1}^{n-1} a\,S_2(n-a)$ with
$S_2(D) = \sum_{d\,\mathrm{odd}\le D} d^{-2} = \pi^2/8 - 1/(2D) +
O(D^{-2})$ (Euler--Maclaurin on the odd-$\zeta(2)$ tail). The tail
contributes $-\sum_{a=1}^{n-1} a/(n-a) = -n H_{n-1} + O(n)$, so this
piece alone gives defect log-coefficient $2$.
\item[(B)] \emph{The $d/2$ correction.} The $(d/2)\cdot(1/d^2)$ piece is
$\sum_{a=1}^{n-1} S_1(n-a)$ with $S_1(D) = \sum_{d\,\mathrm{odd}\le D}
d^{-1} = \tfrac12\log D + O(1)$, which is $\tfrac12 n\log n + O(n)$; it
\emph{adds} to $T_n/Z_n$ and so reduces the defect with log-coefficient
$1$.
\end{itemize}
The remaining pieces --- the $-1/(2a+d)^2$ bracket and the residual
$r(a,d)$ --- carry no $\log n$ at leading order:\ the net coefficient
on the exact $T_n$ is $2 - 1 = 1$. \textbf{[MEASURED]} All three
statements are pinned independently by doubling (Richardson)
estimators on the genuinely computed partial sums, none of which
assumes the answer:\ the piece-(A) estimator converges to $2$, the
piece-(B) estimator to $1$, and the estimator on the exact
$\mathrm{defect}_n$ to $1$
(\texttt{tests/test\_trunk\_qa\_fejer\_4\_over\_pi.py}:\
\texttt{test\_subcoeff\_A\_triangle\_truncation\_is\_two},
\texttt{test\_subcoeff\_B\_sqrt\_correction\_is\_one},
\texttt{test\_net\_defect\_coefficient\_converges\_to\_one}). Hence
\begin{equation*}
   \frac{T_n}{Z_n} \;=\; \frac{\pi^2}{4} \;-\; \frac{\log n}{n}
   \;+\; O\!\Bigl(\frac{1}{n}\Bigr),
\end{equation*}
which is \eqref{eq:Tn_asymptotic}. (An earlier version of this step
attributed the logarithm to the $-1/(2a+d)^2$ bracket and to a boundary
term in $\sqrt{a(a+d)}$, each ``$-\pi^2 n\log n/16$''; that bookkeeping
did not compose to the displayed result and is replaced by the
test-backed decomposition above, 2026-09-02.)

\paragraph{Step 4: Conversion to $\gamma_n$.}
Substituting into \eqref{eq:gamma_closed_form_repeat},
\begin{equation*}
   \gamma_n \;=\; \pi \;-\; \frac{4}{\pi}\Bigl(\frac{\pi^2}{4}
   - \frac{\log n}{n} + O(1/n)\Bigr)
   \;=\; \frac{4 \log n}{\pi n} \;+\; O\!\Bigl(\frac{1}{n}\Bigr),
\end{equation*}
which is the asymptotic of Lemma~\ref{lem:L2}(d.ii).

\paragraph{Numerical verification.}
\textbf{[MEASURED]} The closed form
$n \gamma_n = (4/\pi) \log n + b + O(1/n)$ with
$b \approx 4.109$ has been verified to high precision against the
sum-rule \eqref{eq:gamma_closed_form_repeat} at $n \in \{2, \ldots,
1000\}$, with the doubling estimator $a_n := (2n \gamma_{2n} - n
\gamma_n)/\log 2 \to 4/\pi$ reaching $1.2785$ at $n = 800$ and
$1.2761$ at $n = 1600$ (three digits of $4/\pi$; residual shrinking
${\sim}1.85\times$ per doubling)
(\texttt{tests/test\_trunk\_qa\_fejer\_4\_over\_pi.py}:\
\texttt{test\_gamma\_doubling\_converges\_to\_4\_over\_pi},
\texttt{test\_su2\_doubling\_estimator\_at\_n800} and the slow-marked
\texttt{test\_su2\_doubling\_estimator\_at\_n1600}; the
sum-rule itself is checked against Gaussian quadrature of the kernel
moment in \texttt{tests/test\_central\_fejer\_su2.py}).

\paragraph{Relation to the circle Fej\'er estimate.}
The constant $4/\pi$ is \emph{twice} the constant of the classical
estimate
$\frac{1}{2\pi}\int_{\T^1} F_n(\theta)\,|\theta|\,d\theta \sim (2/\pi)
\log n / n$ for the probability-normalised Fej\'er kernel on the
circle (Zygmund~\cite{zygmund2002} Vol.~I, Ch.~III;
Remark~\ref{rem:circle_fejer}, where the circle constant is computed
in closed form and measured). The $\sin^2(\chi/2)$ factor in the
$\SU(2)$ conjugacy-class measure and the $\sqrt{2j+1}$ Plancherel
weighting of Definition~\ref{def:central_fejer} together change the
kernel and the constant while preserving the leading $\log n / n$
behaviour. \textbf{[OBSERVATION]} On the circle the constant is $2\Vol(S^0)/\Vol(S^1)
= 2/\pi$; on $\SU(2)$ it is $\Vol(S^2)/\Vol(\SU(2)) = 2/\pi$ on the unit
$\sthree$ and $2\Vol(S^2)/\Vol(\SU(2)) = \Vol(S^2)/\pi^2 = 4/\pi$ in the
dual-Coxeter normalisation; no common ``$2\Vol(\text{base})/\Vol(\text{group})$''
form is claimed (2026-09-03), and no derivation of one constant from
the other is given. (An earlier version of this
paragraph stated the circle constant as $4/\pi$ and the $\SU(2)$
constant as ``preserved across this transition''; withdrawn
2026-09-02.)

% =====================================================================
\section*{Acknowledgments}
% =====================================================================

The author thanks Alain Connes, Walter van Suijlekom, Fran\c{c}ois
Latr\'emoli\`ere, Matilde Marcolli, Eva-Maria Hekkelman, Edward
McDonald, Malte Leimbach, and Carlos Perez-Sanchez for the
foundational frameworks on which this paper builds; remaining errors
are the author's own. This work was developed in an AI-augmented
agentic workflow with Anthropic's Claude:\ implementation of the
internal proof-assistant verification machinery, drafting of the
internal proof memos used as supporting documentation, and
bibliographic collation were performed collaboratively. All theorem
statements, design choices, and editorial decisions are the author's
responsibility. The author thanks the open-source NCG community
(notably the SageMath, sympy, and mpmath teams) for the computational
infrastructure used in the supporting numerical verifications.

% =====================================================================
\bibliographystyle{plainnat}
\begin{thebibliography}{99}

\bibitem[Wen \& Avery(1985)]{avery_wen_avery1986}
Z.-Y.~Wen and J.~Avery,
\newblock \emph{Some properties of hyperspherical harmonics},
\newblock J.~Math.~Phys. \textbf{26} (1985), 396--403.

\bibitem[Camporesi \& Higuchi(1996)]{camporesi_higuchi1996}
R.~Camporesi and A.~Higuchi,
\newblock \emph{On the eigenfunctions of the Dirac operator on
spheres and real hyperbolic spaces},
\newblock J.~Geom.~Phys. \textbf{20} (1996), 1--18.

\bibitem[Connes(1995)]{connes1995}
A.~Connes,
\newblock \emph{Noncommutative geometry and reality},
\newblock J.~Math.~Phys. \textbf{36} (1995), 6194--6231.

\bibitem[Connes \& van~Suijlekom(2021)]{connes_vs2021}
A.~Connes and W.~D.~van~Suijlekom,
\newblock \emph{Spectral truncations in noncommutative geometry and
operator systems},
\newblock Comm.~Math.~Phys. \textbf{383} (2021), 2021--2067;
\newblock arXiv:2004.14115.

\bibitem[Hawkins(2000)]{hawkins2000}
E.~Hawkins,
\newblock \emph{Quantization of equivariant vector bundles},
\newblock Comm.~Math.~Phys. \textbf{202} (1999), 517--546;
\newblock and \emph{Geometric quantization of vector bundles and the
correspondence with deformation quantization},
\newblock Comm.~Math.~Phys. \textbf{215} (2000), 409--432.

\bibitem[Hekkelman(2022)]{hekkelman2022}
E.-M.~Hekkelman,
\newblock \emph{Truncated geometry on the circle},
\newblock Lett.~Math.~Phys. \textbf{112} (2022), 20;
\newblock arXiv:2111.13865.

\bibitem[Gaudillot-Estrada \& van~Suijlekom(2025)]{gaudillot_vs2025}
Y.~Gaudillot-Estrada and W.~D.~van~Suijlekom,
\newblock \emph{Convergence of spectral truncations for compact
metric groups},
\newblock Int.~Math.~Res.~Not.\ IMRN \textbf{2025}, no.~13 (2025),
rnaf197;
\newblock arXiv:2310.14733.

\bibitem[van~Suijlekom(2021)]{vs2021_jgp}
W.~D.~van~Suijlekom,
\newblock \emph{Gromov--Hausdorff convergence of state spaces for
spectral truncations},
\newblock J.~Geom.~Phys. \textbf{162} (2021), 104075;
\newblock arXiv:2005.08544.

\bibitem[Rieffel(2004)]{ucp_maps_2024}
M.~A.~Rieffel,
\newblock \emph{Matrix algebras converge to the sphere for quantum
Gromov--Hausdorff distance},
\newblock Mem. Amer. Math. Soc. \textbf{168} (2004), no.~796, 67--91;
arXiv:math/0108005.

\bibitem[Farsi \& Latr\'emoli\`ere(2024)]{farsi_latremoliere2024}
C.~Farsi and F.~Latr\'emoli\`ere,
\newblock \emph{Collapse in noncommutative geometry and spectral
continuity},
\newblock arXiv:2404.00240, 2024.

\bibitem[Farsi \& Latr\'emoli\`ere(2025)]{farsi_latremoliere2025}
C.~Farsi and F.~Latr\'emoli\`ere,
\newblock \emph{Continuity for the spectral propinquity of the
Dirac operators associated with an analytic path of Riemannian
metrics},
\newblock arXiv:2504.11715, 2025.

\bibitem[Hekkelman \& McDonald(2024b)]{hekkelman_mcdonald2024b}
E.~Hekkelman and E.~McDonald,
\newblock \emph{A noncommutative integral on spectrally truncated
spectral triples, and a link with quantum ergodicity},
\newblock J. Funct. Anal., to appear (2025), arXiv:2412.00628.

\bibitem[Toyota(2023)]{toyota2023}
R.~Toyota,
\newblock \emph{Quantum Gromov--Hausdorff convergence of spectral
truncations for groups with polynomial growth},
\newblock arXiv:2309.13469, 2023.

\bibitem[Latr\'emoli\`ere(2015)]{latremoliere2015}
F.~Latr\'emoli\`ere,
\newblock \emph{The dual Gromov--Hausdorff propinquity},
\newblock J.~Math.~Pures Appl. \textbf{103} (2015), 303--351.

\bibitem[Latr\'emoli\`ere(2017)]{latremoliere_metric_st_2017}
F.~Latr\'emoli\`ere,
\newblock \emph{The Gromov--Hausdorff propinquity for metric spectral
triples},
\newblock Adv.~Math. \textbf{404} (2022), Paper No.~108393, 56pp;
\newblock preprint arXiv:1811.10843, 2017.

\bibitem[Latr\'emoli\`ere(2016)]{latremoliere2016}
F.~Latr\'emoli\`ere,
\newblock \emph{The quantum Gromov--Hausdorff propinquity},
\newblock Trans.~Amer.~Math.~Soc. \textbf{368} (2016), 365--411.

\bibitem[Leimbach \& van~Suijlekom(2024)]{leimbach_vs2024}
M.~Leimbach and W.~D.~van~Suijlekom,
\newblock \emph{Gromov--Hausdorff convergence of spectral truncations
for tori},
\newblock Adv.~Math. \textbf{439} (2024), Paper No.~109496.

\bibitem[Marcolli \& van~Suijlekom(2014)]{marcolli_vs2014}
M.~Marcolli and W.~D.~van~Suijlekom,
\newblock \emph{Gauge networks in noncommutative geometry},
\newblock J.~Geom.~Phys. \textbf{75} (2014), 71--91;
\newblock arXiv:1301.3480.

\bibitem[Loutey, internal Paper~32]{paper32}
J.~Loutey,
\newblock \emph{The GeoVac Spectral Triple: A Synthesis Paper Locking
the Framework into the Marcolli--van~Suijlekom Gauge-Network Lineage},
\newblock GeoVac internal preprint (2026), DOI via Zenodo.

\bibitem[Loutey, internal Paper~42]{paper42}
J.~Loutey,
\newblock \emph{Tomita--Takesaki modular structure on truncated
$SU(2)$ spectral triples: four-witness Wick-rotation literal
identification at finite cutoff},
\newblock GeoVac internal preprint (2026), DOI via Zenodo.

\bibitem[Loutey, internal Paper~45]{paper45}
J.~Loutey,
\newblock \emph{Lorentzian propinquity on truncated
$\SU(2) \times U(1)_T$ Krein spectral triples:\ a degeneracy theorem},
\newblock GeoVac internal preprint (2026), DOI via Zenodo.

\bibitem[Loutey, internal Paper~24]{paper24}
J.~Loutey,
\newblock \emph{The Bargmann--Segal Lattice: $\pi$-Free Discretization
of the Harmonic Oscillator on $S^5$},
\newblock GeoVac internal preprint (2026), DOI via Zenodo.

\bibitem[Loutey, internal Paper~40]{paper40_unified}
J.~Loutey,
\newblock \emph{State-space Gromov--Hausdorff convergence of spectral
truncations of compact semisimple Lie groups with bi-invariant metric},
\newblock GeoVac internal preprint (2026), DOI via Zenodo.
\newblock Generalizes the present paper from $\SU(2)$ to compact
connected \emph{semisimple} Lie groups (a circle factor carries the circle
constant); rate constant $4/\pi$ reported universal across that class
(dual-Coxeter normalisation; rigorous at rank~1, rank-uniform proof a
named gap) via a Plancherel-weight $\times$ Vandermonde-Jacobian
cancellation; asymptotic-tight $C_3 = 1$ at all ranks via the PRV-summand
bound (Kumar 1988) with a Brauer--Klimyk interior closure.

\bibitem[Loutey, internal Paper~2]{paper2}
J.~Loutey,
\newblock \emph{The Fine Structure Constant as a Spectral Coincidence
on the Hopf Fibration},
\newblock GeoVac internal preprint, papers/observations (2026),
DOI via Zenodo.

\bibitem[Loutey, internal Paper~7]{paper7}
J.~Loutey,
\newblock \emph{The Dimensionless Vacuum: Recovering the Schr\"odinger
Equation from Scale-Invariant Graph Topology},
\newblock GeoVac internal preprint (2024), DOI via Zenodo.

\bibitem[Loutey, internal Paper~18]{paper18}
J.~Loutey,
\newblock \emph{Spectral-Geometric Exchange Constants: Projecting
Discrete Spectra onto Continuous Manifolds. Master Mellin engine
\S III.7},
\newblock GeoVac internal preprint, papers/group3\_foundations (2026),
DOI via Zenodo.

\bibitem[Perez-Sanchez(2025)]{perez_sanchez2025}
C.~I.~Perez-Sanchez,
\newblock \emph{Comment on `Gauge networks in noncommutative
geometry'},
\newblock arXiv:2508.17338, 2025.

\bibitem[Stein \& Weiss(1971)]{stein_weiss1971}
E.~M.~Stein and G.~Weiss,
\newblock \emph{Introduction to Fourier Analysis on Euclidean Spaces},
\newblock Princeton Math.~Ser. \textbf{32}, Princeton Univ.~Press, 1971.

\bibitem[Zygmund(2002)]{zygmund2002}
A.~Zygmund,
\newblock \emph{Trigonometric Series, Vol.~I and II},
\newblock 3rd ed., Cambridge Univ.~Press, 2002.

\bibitem[Loutey, internal Paper~43]{paper43}
J.~Loutey,
\newblock \emph{Lorentzian extension of the four-witness Wick-rotation
theorem on truncated $\sthree \times \R$ spectral triples at finite
cutoff},
\newblock GeoVac internal preprint (2026), DOI via Zenodo.

\end{thebibliography}

\end{document}
